% Options for packages loaded elsewhere
\PassOptionsToPackage{unicode}{hyperref}
\PassOptionsToPackage{hyphens}{url}
%
\documentclass[
  10pt,
]{article}
\usepackage{amsmath,amssymb}
\usepackage{iftex}
\ifPDFTeX
  \usepackage[T1]{fontenc}
  \usepackage[utf8]{inputenc}
  \usepackage{textcomp} % provide euro and other symbols
\else % if luatex or xetex
  \usepackage{unicode-math} % this also loads fontspec
  \defaultfontfeatures{Scale=MatchLowercase}
  \defaultfontfeatures[\rmfamily]{Ligatures=TeX,Scale=1}
\fi
\usepackage{lmodern}
\ifPDFTeX\else
  % xetex/luatex font selection
\fi
% Use upquote if available, for straight quotes in verbatim environments
\IfFileExists{upquote.sty}{\usepackage{upquote}}{}
\IfFileExists{microtype.sty}{% use microtype if available
  \usepackage[]{microtype}
  \UseMicrotypeSet[protrusion]{basicmath} % disable protrusion for tt fonts
}{}
\makeatletter
\@ifundefined{KOMAClassName}{% if non-KOMA class
  \IfFileExists{parskip.sty}{%
    \usepackage{parskip}
  }{% else
    \setlength{\parindent}{0pt}
    \setlength{\parskip}{6pt plus 2pt minus 1pt}}
}{% if KOMA class
  \KOMAoptions{parskip=half}}
\makeatother
\usepackage{xcolor}
\usepackage[margin=0.85in]{geometry}
\setlength{\emergencystretch}{3em} % prevent overfull lines
\providecommand{\tightlist}{%
  \setlength{\itemsep}{0pt}\setlength{\parskip}{0pt}}
\setcounter{secnumdepth}{-\maxdimen} % remove section numbering
\DeclareUnicodeCharacter{03B8}{\ensuremath{\theta}}
\DeclareUnicodeCharacter{03B5}{\ensuremath{\epsilon}}
\DeclareUnicodeCharacter{03C6}{\ensuremath{\varphi}}
\DeclareUnicodeCharacter{03BC}{\ensuremath{\mu}}
\DeclareUnicodeCharacter{03BD}{\ensuremath{\nu}}
\DeclareUnicodeCharacter{03B1}{\ensuremath{\alpha}}
\DeclareUnicodeCharacter{03B4}{\ensuremath{\delta}}
\DeclareUnicodeCharacter{03C9}{\ensuremath{\omega}}
\DeclareUnicodeCharacter{03C1}{\ensuremath{\rho}}
\DeclareUnicodeCharacter{2265}{\ensuremath{\geq}}
\DeclareUnicodeCharacter{2264}{\ensuremath{\leq}}
\DeclareUnicodeCharacter{2212}{\ensuremath{-}}
\DeclareUnicodeCharacter{2208}{\ensuremath{\in}}
\DeclareUnicodeCharacter{2229}{\ensuremath{\cap}}
\DeclareUnicodeCharacter{2192}{\ensuremath{\to}}
\DeclareUnicodeCharacter{221E}{\ensuremath{\infty}}
\DeclareUnicodeCharacter{222B}{\ensuremath{\int}}
\DeclareUnicodeCharacter{220F}{\ensuremath{\prod}}
\DeclareUnicodeCharacter{03A3}{\ensuremath{\sum}}
\DeclareUnicodeCharacter{2225}{\ensuremath{\Vert}}
\DeclareUnicodeCharacter{2272}{\ensuremath{\lesssim}}
\DeclareUnicodeCharacter{00D7}{\ensuremath{\times}}
\DeclareUnicodeCharacter{2227}{\ensuremath{\wedge}}
\DeclareUnicodeCharacter{2282}{\ensuremath{\subset}}
\DeclareUnicodeCharacter{2297}{\ensuremath{\otimes}}
\DeclareUnicodeCharacter{2205}{\ensuremath{\varnothing}}
\DeclareUnicodeCharacter{2194}{\ensuremath{\leftrightarrow}}
\ifLuaTeX
  \usepackage{selnolig}  % disable illegal ligatures
\fi
\usepackage{bookmark}
\IfFileExists{xurl.sty}{\usepackage{xurl}}{} % add URL line breaks if available
\urlstyle{same}
\hypersetup{
  pdftitle={A relative finite-energy equation on the entire big rational surface cone},
  pdfauthor={Unreviewed mathematical candidate},
  hidelinks,
  pdfcreator={LaTeX via pandoc}}

\title{A relative finite-energy equation on the entire big rational
surface cone}
\author{Unreviewed mathematical candidate}
\date{6 October 2026}

\DeclareUnicodeCharacter{1D4E7}{\ensuremath{\mathcal X}}
\DeclareUnicodeCharacter{00B9}{\ensuremath{^1}}
\begin{document}
\maketitle

\section{A relative equation on the entire big rational surface
cone}\label{a-relative-equation-on-the-entire-big-rational-surface-cone}

This is an unreviewed candidate. The theorem below concerns every actual
big rational bundle on a smooth integral projective surface over every
algebraically closed trivially valued characteristic-zero field. It does
not claim the arbitrary-dimensional original P16.5. The complete surface
structural proof and the complete proposed intrinsic nef operator,
polar, separation and attainment proofs are retained in the appendices.
None of their historical status labels is a grant for the present
theorem. The only analytic grant is the exact whole actual-ample \(G_0\)
body in its retained domain; the new combined argument requires
independent verification.

\subsection{S1. Retained datum and the claimed
equation}\label{s1.-retained-datum-and-the-claimed-equation}

Let \(k\) be as above, \(X\) a smooth integral projective surface, \(L\)
any actual big rational bundle and \(H\) an actual ample rational
bundle. There is no countability, semiampleness, finite-generation,
toric, skeletal or support restriction. Appendix A constructs the
canonical actual decomposition

\[L=P+N,\qquad P\text{ nef},\quad N\geq0,\quad P\cdot C=0\ (C\subset\operatorname{Supp}N),\quad V=P^2=\operatorname{vol}(L)>0.\]

It proves, at every original analytic point and for arbitrary directed
nets,

\[\operatorname{PSH}(L)=\sigma_N+\operatorname{PSH}(P),\qquad V_L=\sigma_N.\]

Here the left cone is precisely one all-point decreasing net of actual
generically finite FS section maxima with numerical classes tending to
\(c_1(L)\) and finite generic-trivial endpoint. The entire normalized
cone is divisorially compact and decreasingly closed. The divisor
potential \(\sigma_N\) uses the actual canonical section, and every
remainder \(v\) is an actual uniquely determined function on the
original \(X^{\mathrm{an}}\). Divisorial subtraction is finite; equality
everywhere follows from common-ample determination. The compensator is
\(\mathcal C_L(\sigma_N+v)=v\), including its specified values at shared
poles.

\textbf{Candidate theorem.} Define the relative energy and entire
finite-energy space by

\[E_L(\sigma_N+v)=E_P(v),\qquad \mathcal E_L^1=\sigma_N+\mathcal E^1(P).\]

There is an actual target-independent moving-intersection operator
\(M_L\) on this whole domain, of raw mass \(V\). It has comparison,
domination, decreasing mixed and compensator-weighted continuity, both
outer derivatives in every original continuous direction, and
all-solution uniqueness modulo constants. For every positive Radon
measure \(\nu\) on \(X^{\mathrm{an}}\) of mass \(V\), the equation
\(M_L(u)=\nu\) has an actual solution if and only if

\[D_L(\nu):=\sup_{u\in\mathcal E_L^1}\left(E_L(u)-\int\mathcal C_L(u)\,d\nu\right)<\infty.\]

Such a solution actually attains the displayed supremum. This is the
intrinsic compensator condition, with raw mass and reference specified
before the target. It will itself imply polar exclusion; there is no
extra support hypothesis.

\subsection{S2. The whole intrinsic nef theorem used, with its proof
included}\label{s2.-the-whole-intrinsic-nef-theorem-used-with-its-proof-included}

For the actual nef \(P\), let \(v_j=\max(v,-j)\) after generic
normalization. The bounded products are the genuine signed
metric-pairing products in a common actual ample finite-energy domain,
with positive ample perturbations proving positivity and raw mass
\(P^2\). They retain the model central-component multiplicities.
Appendices B--E reconstruct

\[E_P(v)=\inf_j E_P(v_j),\qquad \mathcal E^1(P)=\{v\in\operatorname{PSH}(P):E_P(v)>-\infty\},\]

and the target-independent product from the increasing strict-winning
restrictions of the actual bounded floor measures. The floor
total-variation error is at most \(2(n+1)(-E_P(v))/j\) for normalized
\(v\), here \(n=2\). This retains the entire intrinsic domain, which can
exceed the auxiliary ample pairing domain. No unbounded intrinsic
endpoint is simply placed in that auxiliary space.

The appended proofs establish every common-class mixed product, its full
mass, integrability of every intrinsic finite-energy test, arbitrary
decreasing unbounded and weighted convergence, both continuous-direction
outer derivatives, weak energy upper semicontinuity, compact cost
sublevels, intrinsic dual finiteness equivalence and coercivity. The
additional polar, finite-divisorial atomic and PL approximation
arguments prove all-point separation of \(I_P\), target-integral
continuity on every cost sublevel, and actual attainment for every
finite-dual-energy target in the positive big nef class. The target
integration argument uses finite divisorial model measures and uniform
numerical estimates; energy upper semicontinuity alone is insufficient.

These are submitted proofs, rather than a grant obtained from a
historical PASS label. In particular, the earlier review125 had an
invalid evidence index and did not produce an accepted theorem. All its
necessary mathematical mechanisms are included here for fresh
verification. The appendices retain their exact earlier mathematical
scope statements; their local equation numbers belong to the individual
appendix.

\subsection{S3. A genuine pole test for the negative
divisor}\label{s3.-a-genuine-pole-test-for-the-negative-divisor}

Put \(Z=(\operatorname{Supp}N)^{\mathrm{an}}\). We prove that every
intrinsic finite-dual-energy \(P\) target, and every actual \(P\)
output, gives \(Z\) zero mass. If \(N=0\) this is immediate.

Choose a sufficiently small rational \(\delta>0\) so that
\(B=P-\delta H\) has \(B^2>0\) and \(B\cdot H>0\). For large divisible
\(m\), \(K_X-mB\) has negative \(H\) degree, so Serre duality and
surface Riemann--Roch give

\[h^0(mB)\geq\chi(mB)=\tfrac12m^2B^2+O(m)>0.\]

An actual nonzero section gives an effective rational divisor \(D\) with
\(B=\mathcal O(D)\) in the rational Picard group. Thus \(P=\delta H+D\)
as actual rational bundles. Let \(G=\lceil N\rceil\), an effective
integral divisor. For a sufficiently large divisible \(q\), the actual
ample \(q\delta H-G\) has a nonzero section. Multiplying it by the
canonical sections of \(G\) and \(qD\) gives a nonzero section
\(h\in H^0(qP)\) vanishing on \(\operatorname{Supp}N\).

Its normalized one-section potential \(\tau=q^{-1}\log|h|\) is genuine
\(P\)-psh and is \(-\infty\) on \(Z\). The bounded floor
\(b_T=\max(\tau,-T)\) is genuine \(P\)-psh, generic-normalized and
between \(-T\) and zero. The intrinsic energy therefore satisfies
\(0\leq -E_P(b_T)/V\leq T\). The reconstructed quadratic contraction
estimate gives

\[-E_P(a b_T)/V\leq K_2a^2T,\qquad 0<a\leq1,\]

where \(K_2\) is a fixed dimension constant. If \(m=\nu(Z)/V>0\), then

\[\frac1V\left(E_P(a b_T)-\int a b_T\,d\nu\right)\geq(am-K_2a^2)T.\]

Choose a fixed sufficiently small rational \(a>0\). The right side tends
to infinity. Thus finite intrinsic dual energy forces \(\nu(Z)=0\).
Every top output has finite dual by Appendix B. Every common-class mixed
output does as well, by its proved intrinsic-test integrability and the
dual criterion. Hence all these outputs vanish on \(Z\). This argument
uses an actual big nef section and full vanishing, with no countable
decomposition of polar sets and no assumption that \(\sigma_N\) itself
is \(P\)-psh.

\subsection{S4. Construct the relative product before the target
equation}\label{s4.-construct-the-relative-product-before-the-target-equation}

On \(U=X\setminus\operatorname{Supp}N\), the actual canonical \(N\)
section is invertible. Removing it in the frame of \(L\) identifies the
metric of \(u=\sigma_N+v\) with the metric of the actual bundle \(P\)
and potential \(v\). At every finite nef model stage, removing this
fixed horizontal factor leaves exactly the nef metric line bundle for
\(P\). Its mixed central-component intersection coefficients are
therefore exactly those of the genuine bounded \(P\) product. They
include the same central multiplicities. This constructs the local
moving product from actual model intersections, independent of any
prescribed target.

For bounded genuine remainders, use the bounded metric-pairing
construction of Appendix B and its actual positive ample perturbations;
finite model refinements identify the same local moving product on
\(U^{\mathrm{an}}\). The product has zero mass on \(Z\) by S3, so its
nonpluripolar extension to the original compact carrier is uniquely
specified by zero polar mass. Define

\[M_L(\sigma_N+v)=M_P(v).\]

This identity records an already constructed moving product. It does not
select an inverse measure in response to a target. For every unbounded
\(v\in\mathcal E^1(P)\), construct the relative floors \(\sigma_N+v_j\).
Their actual moving products have the strict-winning and total-variation
limits proved in Appendix B. These define the displayed original-carrier
operator on the entire relative domain, with raw mass \(V\) and no
auxiliary ample dependence.

The same construction in every mixed slot gives the common-\(L\)
relative mixed products. For decreasing genuine \(u_i=\sigma_N+v_i\),
the remainders decrease all-pointwise: divisorial decrease is finite,
and common-ample genuine maxima promote it to all points. Therefore all
raw energy, mixed and weighted limits follow from the appended intrinsic
nef proofs. The weights are specifically \(\mathcal C_L(w_i)\), not bare
logarithmic potentials including the fixed divisor. Since the weighted
tests are integrable and the underlying outputs vanish on \(Z\), their
signed weighted measures have zero mass on \(Z\) as well. All limits are
on the original full continuous test space.

\subsection{S5. Genuine envelopes and both outer
derivatives}\label{s5.-genuine-envelopes-and-both-outer-derivatives}

We first retain the precise order fact: if \(a,b\) are already genuine
and \(f\) is continuous, \(a\leq b+f\) at all divisorial points implies
it at every analytic point. Approximate \(f\) uniformly by real PL
functions, express each as a difference of bounded genuine ample
potentials, put both compared sums in a common sufficiently large actual
ample cone by zero-system tensoring, and use ample divisorial
order/determination. Let the uniform error decrease to zero. The same
argument compares a genuine potential with a continuous obstacle.

For any upper-bounded continuous obstacle, constants below it added to
\(\sigma_N\) supply genuine minorants. After replacing competitors by
their maximum with a fixed lower minorant, their generic constants lie
in a compact interval. Normalized whole-cone compactness and the closed
divisorial constraints then give a greatest actual minorant: finite
maxima have a cluster point with the supremum coordinates, and the
preceding order fact makes it greatest all-pointwise. Thus \(P_L\) is a
genuine envelope, without importing a general big envelope property.

Now fix \(u=\sigma_N+v\in\mathcal E_L^1\) and
\(f\in C^0(X^{\mathrm{an}})\). At every divisor \(\sigma_N\) is finite.
An actual competitor \(\sigma_N+w\) is below \(u+tf\) precisely when
\(w\leq v+tf\) at all divisors, and then everywhere by the order fact.
The latter implication also follows on \(U\) by literal finite
subtraction, while at \(Z\) the original inequality is automatic.
Consequently

\[P_L(u+tf)=\sigma_N+P_P(v+tf).\]

The right side is in the entire relative finite-energy domain, between
\(u-|t|\|f\|\) and \(u+|t|\|f\|\). For decreasing continuous majorants
of \(u+tf\), the greatest genuine \(L\) minorants decrease to this same
envelope by proved genuine closure and all-point constraints. Their
remainders decrease; intrinsic nef energy continuity identifies the
infimum defining the outer energy with \(E_L(P_L(u+tf))\). The full
original continuous-direction derivative of the nef proof now gives both
signs:

\[\left.\frac d{dt}\right|_{0^+}E_L^\downarrow(u+tf)=\left.\frac d{dt}\right|_{0^-}E_L^\downarrow(u+tf)=\int f\,M_P(v)=\int f\,M_L(u).\]

The test function stays on the original \(X\) throughout. Neither its
support near \(Z\) nor its values there have been removed as a
condition; the intrinsic operator itself gives that polar set zero mass.

\subsection{S6. Exact targets, attainment and
uniqueness}\label{s6.-exact-targets-attainment-and-uniqueness}

By the entire cone and compensator bijections, for every positive
raw-mass-\(V\) Radon target,

\[D_L(\nu)=\sup_{v\in\mathcal E^1(P)}\left(E_P(v)-\int v\,d\nu\right)=D_P(\nu).\]

Constants cancel by the same raw mass. Finiteness is equivalent to
integrability of every intrinsic compensator by the appended nef
criterion, and it implies \(\nu(Z)=0\) by S3. The full intrinsic nef
variational proof gives an actual \(v\in\mathcal E^1(P)\) attaining the
supremum and satisfying \(M_P(v)=\nu\). The single actual FS
construction of Appendix A then gives \(u=\sigma_N+v\) as an actual
function on the original carrier. By S4 its preconstructed moving
operator satisfies \(M_L(u)=\nu\), and the relative objective is the
same attained objective. Conversely every solution has finite intrinsic
dual because each actual nef output does. This proves the exact range
criterion without an independent support restriction.

For all-solution uniqueness, equality of the relative operators gives
equality of the actual nef operators. The appended finite-divisorial
atomic equations and four-point numerical inequality prove separation of
\(I_P\) at every divisor. Common-ample determination gives \(v=w+C\) at
all points, hence \(u=u'+C\) in the metric sense, including shared
poles. No undefined difference of infinities is used.

\subsection{S7. Comparison and domination in the whole
domain}\label{s7.-comparison-and-domination-in-the-whole-domain}

For bounded genuine \(P\) potentials, raw comparison passes from actual
ample \(P+\epsilon H\) through the finite measure-polynomial
coefficients: at the fixed Borel winning set, an inequality for every
positive \(\epsilon\) retains its constant coefficient as
\(\epsilon\downarrow0\). For intrinsic finite-energy inputs, literal
floors converge in total variation. The bounded winning set for
\(v_j<w_j\) is \(\{v<w\}\cap\{w>-j\}\) and increases to \(\{v<w\}\)
except on common poles, which all actual outputs give zero mass.
Ordinary countable convergence and total variation therefore give

\[\int_{\{v<w\}}M_P(w)\leq\int_{\{v<w\}}M_P(v).\]

The same floor argument retains strict-winning locality. On
\(U^{\mathrm{an}}\), adding the same finite \(\sigma_N\) preserves the
winning sets. All relative measures vanish on \(Z\). Thus comparison and
locality hold for all actual relative finite-energy potentials.

Domination also follows without an unproved nonample source theorem. If
\(v\leq w\) almost everywhere for \(M_P(w)\), put \(r=\max(v,w)\). The
supporting inequality and monotonicity give

\[0\leq E_P(r)-E_P(w)\leq\int(r-w)M_P(w)=0.\]

The ordered numerical estimate forces \(I_P(r,w)=0\). Its proved
all-point separation gives \(r=w+C\); the zero energy difference and the
translation mass \(V>0\) force \(C=0\). Hence \(v\leq w\) everywhere.
The relative almost-everywhere condition subtracts the finite fixed
factor on \(U\); polar exclusion makes it exactly the same hypothesis
for the nef remainders. Their all-point order then yields the relative
order on the full carrier. Positive volume is used in the numerical
separation; no zero-mass domination is claimed.

\subsection{S8. Additional zero-volume surface range and the exact
original
boundary}\label{s8.-additional-zero-volume-surface-range-and-the-exact-original-boundary}

Appendix A also treats every rational pseudoeffective surface class,
including classes with no sections. When its canonical positive part is
numerically zero, its entire genuine cone is \(\sigma_N+\mathbb R\). The
actual relative moving products and energy are zero; the only positive
target is zero and solutions are unique modulo constants by cone
rigidity. This is an explicit additional zero-volume range, with no
probability normalization. It includes nonzero negative exceptional
classes.

For other zero-volume positive parts, the zero top operator need not be
injective: bounded normalized potentials on \(\mathcal O(1,0)\) on
\(\mathbb P^1\times\mathbb P^1\) already have identical zero products
and energies while differing nonconstantly. The appended construction
retains this obstruction. No unrestricted zero-volume uniqueness theorem
is inferred.

The candidate therefore combines an entire positive-volume big-surface
equation with this precise additional rigid zero-volume range. It does
not prove arbitrary-dimensional nonnef big closure or equation, nor a
unique positive-mass replacement for all nonzero zero-volume classes.
Those remain part of the full unchanged original P16.5, which is open.
No original or accepted-paper count changes until an exact independent
review has been read and validated.

\section{Appendix A. Complete pseudoeffective surface structure and
exact zero
ranges}\label{appendix-a.-complete-pseudoeffective-surface-structure-and-exact-zero-ranges}

The following complete mathematical derivation is retained as supporting
proof, with its earlier scope statements. It is not a historical-verdict
grant. Equation numbers in this appendix are local to it. The main
candidate and the fresh review criteria govern the new combined scope.

\section{1. Genuine objects, retained hypotheses, and the premise
boundary}\label{genuine-objects-retained-hypotheses-and-the-premise-boundary}

Unless a dimension is specified, let \(k\) be an arbitrary algebraically
closed field of characteristic zero with the trivial absolute value, and
\(X\) a smooth integral projective \(k\)-variety of positive dimension.
There is no cardinality restriction on \(k\). An actual rational line
bundle means an element of \(\operatorname{Pic}(X)\otimes\mathbb Q\);
every displayed section uses an integral tensor power clearing its
denominator. Fix an actual ample integral line bundle \(H\), with
\(h=c_1(H)\). Taking a sufficiently divisible power allows this choice
without restricting the conclusions. Write \(x_0\) for the generic
trivial valuation. Every potential below is an actual function on the
original compact Hausdorff Berkovich space \(X^{\rm an}\).

For an actual rational bundle \(B\), put \[
 \operatorname{FS}^{\rm gf}(B)=\left\{
  \frac1m\max_j(\log\|s_j\|_{\rm triv}+\lambda_j):
 0\ne s_j\in H^0(X,mB),\ \lambda_j\in\mathbb R
 \right\}.                                                     \tag{1.1}
\] There are finitely many sections in each formula. Their base ideal
can be nontrivial and their maximum can equal \(-\infty\). The canonical
trivial metric of the specified actual bundle is used. We abbreviate
\(\log\|s\|_{\rm triv}\) to \(\log|s|\). The subfamily
\(\operatorname{FS}^{\rm fin}(B)\) consists of functions finite at every
analytic point. These are continuous real-valued functions. An
extended-valued section maximum is continuous into
\(\mathbb R\cup\{-\infty\}\).

For \(\theta\in N^1(X)_{\mathbb R}\) define the genuine \textbf{one-net}
cone by \[
 \operatorname{PSH}(\theta)=\left\{u:
 u(x_0)>-\infty,\quad u=\inf_i f_i\text{ at every analytic point},
 \quad f_i\searrow,\quad f_i\in\operatorname{FS}^{\rm gf}(B_i),
 \quad c_1(B_i)\longrightarrow\theta\right\}.                    \tag{1.2}
\] Every index set here is an arbitrary directed set. No abstract
completion, transfinite closure, or weak representative replaces (1.2).
An infimum of these continuous extended-valued functions is usc. At a
valuation-ring center all regular section norms are at most one, while a
nonzero section has norm one at \(x_0\). Consequently \[
 \sup_{X^{\rm an}} f=f(x_0)=m^{-1}\max_j\lambda_j,
 \qquad \sup u=u(x_0).                                          \tag{1.3}
\] The normalized cone has \(u(x_0)=0\). Its topology is pointwise
convergence at divisorial valuations of \(k(X)/k\). An arbitrary
positive real multiple of a divisorial valuation is permitted. This
topology is not all-point convergence.

The only analytic grant, denoted \(G_0\), is the exact supplied whole
P16.1 body \textbf{in its actual ample rational-line-bundle domain}. On
our smooth original \(X,k\) we use the following actual-ample
consequences: finite continuous same-bundle decreasing FS
regularization; decreasing closure with a finite generic limit;
compactness of the normalized cone in divisorial coordinates; finite
divisorial evaluations and all-point determination by those coordinates;
finite maxima; and the raw actual-ample model intersection product. No
nonample consequence mentioned in the commentary of \(G_0\) is a
premise. Its equation theorem is not applied to a nef bundle merely by
changing the adjective ample. The supplied completed132 paper, root183
candidate, and completed129 report are mathematical data to audit, not
additional grants.

The ordinary algebraic background is explicitly retained: finite
dimensionality of \(N^1(X)_{\mathbb R}\); numerical invariance and
openness of ampleness; generation of sufficiently divisible ample
powers; closedness of the nef cone; nef plus ample is ample; positivity
of distinct-curve intersections and of \(C\cdot H\) on a smooth surface;
the finite polytope and real symmetric-matrix spectral theorems; Serre
vanishing, Serre duality, and surface Riemann--Roch. For the multiplier
calculation we use the definition by a log resolution and the elementary
regular-sequence associated graded algebra. These standard algebraic
theorems are stated prerequisites, not analytically granted nonample
results. All Zariski decomposition, negative-part continuity, section
division, and nonample cone arguments used here are proved below. No
Nadel or Kawamata--Viehweg vanishing theorem is invoked.

The supplied literal primary unit is BJ, \emph{Global pluripotential
theory over a trivially valued field}, arXiv:1801.08229v4, original
LF2180--2335. Its definition motivates (1.2). Its extended-real
flattening lemma is not used. The other reported primary coordinates in
the historical documents are not represented as freshly read source
units. A public request for Bauer, arXiv:0712.1576v1, failed at the
configured proxy; no Bauer proof was read in this workspace. Our finite
decomposition proof is written in full.

\section{2. Actual section operations and the repaired nef
foundation}\label{actual-section-operations-and-the-repaired-nef-foundation}

\textbf{Lemma 2.1 (actual zero systems).} An actual ample rational
bundle \(C\) admits a generated power whose weight-zero FS potential is
identically zero. Tensor products realize sums of FS potentials,
retaining actual base ideals and real weights.

\textbf{Proof.} Choose generators \(t_b\) of a divisible \(qC\). At each
valuation-ring center one generator is a unit of norm one and all
generators have norm at most one. Hence \(q^{-1}\max_b\log|t_b|=0\) at
every analytic point, including semivaluations with nonzero kernel. To
tensor two formulas, pass to a common multiple \(M\) of their degrees.
Use all products of \(M/m\) sections of the first system and all
products of \(M/q\) of the second, with summed weights. Their maximum is
exactly the sum of the two original maxima, also at poles. Locally their
base ideal is the product of the corresponding powered actual ideals. In
particular tensoring a zero system embeds a function in
\(\operatorname{FS}^{\rm gf}(B)\) unchanged into
\(\operatorname{FS}^{\rm gf}(B+C)\). \(\square\)

\textbf{Lemma 2.2 (one-net class embedding).} If
\(u\in\operatorname{PSH}(\theta)\) and \(T\) is an actual rational
bundle with \(c_1(T)-\theta\) ample, then the same all-point function
belongs to \(\operatorname{PSH}(c_1(T))\). The target \(T\) itself need
not be ample.

\textbf{Proof.} On a tail of its original defining net, \(T-B_i\) is
actual ample by numerical ampleness and openness. Tensor the zero
systems from Lemma 2.1 separately at every stage. The unchanged \(f_i\)
are now actual FS functions on the same \(T\), with the original
all-point order and limit. This is already one defining net. \(\square\)

In particular any finitely many already genuine functions with specified
classes can be put into one sufficiently large actual ample cone. They
then have finite divisorial coordinates. Equality of those coordinates
implies all-point equality by \(G_0\). This conclusion requires
\textbf{prior genuine membership} of both functions. It does not prove
membership for a new limit.

\textbf{Lemma 2.3 (nef intersection, with a single net).} For every
actual nef rational \(P\), \[
 \operatorname{PSH}(c_1(P))
   =\bigcap_{\epsilon\in\mathbb Q_{>0}}
                      \operatorname{PSH}(P+\epsilon H).            \tag{2.1}
\] Every function on the left admits a single decreasing all-point net
in \(\operatorname{FS}^{\rm fin}(P+\epsilon_aH)\) with
\(\epsilon_a\to0\).

\textbf{Proof.} Forward inclusion is Lemma 2.2. Conversely let \(u\) be
in every actual ample cone on the right. For every rational
\(\epsilon>0\) choose its \(G_0\) finite continuous regularizers
\(f_{\epsilon,j}\searrow u\) on the actual \(P+\epsilon H\). Index
triples \((\epsilon,j,b)\) with \(b>0\) by simultaneous decrease of
\(\epsilon\), of \(b\), and of the function \(f_{\epsilon,j}+b\).

This preorder is directed. Given two triples, choose \(\epsilon_3\)
below their \(\epsilon\) values and \(0<b_3<\min(b_1,b_2)\). The
function \[
 R=\min(f_{\epsilon_1,j_1}+b_1,f_{\epsilon_2,j_2}+b_2)
\] is finite continuous. At every point \(u+b_3<R\), including where
\(u=-\infty\). Thus the open negative sets of \[
 \exp(f_{\epsilon_3,j}+b_3)-\exp R
\] cover the compact carrier. They increase with \(j\). A finite
subcover and directedness give one \(j_3\) with strict comparison
everywhere. This triple is later than both original triples.

The resulting shifted FS net decreases and stays above \(u\). For any
point \(x\) and any finite threshold \(a>u(x)\), choose one \(j,b\) with
\(f_{\epsilon,j}(x)+b<a\). Every later triple stays below that
threshold. This treats \(u(x)=-\infty\) without a meaningless strict
comparison between two infinities. For every \(\delta>0\) a triple with
\(\epsilon<\delta\) exists, and all later triples have that property.
The numerical classes tend to \(c_1(P)\); the limit is exactly \(u\) at
all points. The generic value is finite since \(u\) belongs to an ample
cone. This proves (1.2) and (2.1). \(\square\)

\textbf{Proposition 2.4 (entire nef closure and compactness).} The
entire genuine nef rational cone is closed under arbitrary decreasing
nets with finite limiting generic value. Its normalized part is compact
in the original divisorial topology.

\textbf{Proof.} A decreasing net in the nef cone lies, unchanged, in
each actual ample \(P+\epsilon H\). Ample decreasing closure gives its
all-point limit there. Lemma 2.3 supplies one defining net in the nef
class. For compactness, let \(K_\epsilon\) be the normalized actual
ample cone for \(P+\epsilon H\), \(0<\epsilon\leq1\). Each is compact by
\(G_0\) and embeds in the Hausdorff coordinate space \(K_1\). Its
compact image is closed. Their intersection is compact. If its
description initially gives a representative only through divisorial
coordinates, the representatives all belong to the same ample \(K_1\),
so all-point determination identifies them. Equation (2.1) then
identifies the intersection with the entire genuine normalized nef cone,
with its actual functions. \(\square\)

\textbf{Lemma 2.5 (varying nef compact limit).} Suppose \(P_a\) are
actual nef rational bundles, \(c_1(P_a)\to c_1(P)\) for an actual nef
rational \(P\), and \(v_a\in\operatorname{PSH}(c_1(P_a))\) are
normalized. A subnet converges in a common normalized actual ample cone;
any such limit is a genuine normalized \(P\)-psh function, with one
defining net.

\textbf{Proof.} Choose actual ample \(A\) so large that \(A-P-H\) is
ample. On a tail \(A-P_a\) is ample, so Lemma 2.2 puts the same \(v_a\)
into compact \(K_A\). Choose a convergent subnet. For each rational
\(0<\delta\leq1\), eventually \(P+\delta H-P_a\) is ample. The \(v_a\)
then belong to the compact normalized actual ample cone
\(K_{P+\delta H}\), whose image in \(K_A\) is closed. The actual
representative of the limit lies in every such cone, by ample
determination. Lemma 2.3 gives the required genuine \(P\)-psh function
and its one net. This uses subnets, not countable extraction.
\(\square\)

These proofs reconstruct the necessary completed132 nef results. The
finite-valued nature of the approximants in Lemma 2.3 is indispensable.
At a shared pole, a strict inequality between two extended-real
approximants is false. The compact open-cover argument here works
because \(R\) is finite.

\section{3. Complete finite algebraic decomposition on a
surface}\label{complete-finite-algebraic-decomposition-on-a-surface}

From now through Section 7, \(\dim X=2\). A bundle is pseudoeffective if
its numerical class lies in the closure of the cone of effective
rational divisors. An actual big bundle has positive section volume; in
particular it has an effective rational divisor representative. We first
work with such a representative, then remove the existence-of-sections
restriction.

\textbf{Lemma 3.1 (maximal nef subdivisor).} An effective rational
divisor \(D\) has a decomposition \[
 D=P_D+N_D,
 \quad P_D\text{ effective nef},\quad N_D\geq0,\quad
 P_D\cdot C=0\text{ on }\operatorname{Supp}N_D,
                                                               \tag{3.1}
\] whose negative support has negative definite intersection matrix.
Both parts are rational.

\textbf{Proof.} Write \(D=\sum_{i=1}^r d_iC_i\) with distinct prime
curves. The compact rational polytope \[
 K=\left\{(p_i):0\leq p_i\leq d_i,\quad
                  \sum_i p_i(C_i\cdot C_j)\geq0\ (1\leq j\leq r)
      \right\}                                                \tag{3.2}
\] is nonempty. Its points are exactly the effective nef subdivisors:
intersections with every curve outside this finite support are
nonnegative. Maximize \(\sum p_i\) and choose a rational maximizing
vertex.

Coordinatewise maxima preserve feasibility. For the inequality tested
against \(C_j\), choose the original feasible vector realizing the
larger \(j\)th coordinate. The diagonal contribution stays fixed and all
other coordinates increase with nonnegative off-diagonal coefficients.
Consequently the maximizing vector is the unique greatest feasible
vector: a feasible vector exceeding it in any coordinate would make
their coordinatewise maximum have larger sum. Set \(P_D=\sum p_iC_i\),
\(N_D=D-P_D\).

If \(a_j=d_j-p_j>0\) and \(P_D\cdot C_j>0\), a small increase of \(p_j\)
preserves the \(j\)th inequality and the cuboid; all other tested
intersections only increase. This contradicts maximality. Thus the
claimed orthogonality holds.

Let \(Q\) be the intersection matrix on the support of \(N_D\). If its
largest eigenvalue were nonnegative, choose a unit vector maximizing its
Rayleigh quotient. Replacing it by its coordinatewise absolute value
cannot decrease that quotient, since the off-diagonal entries are
nonnegative. A nonnegative maximizing vector \(q\ne0\) is an
eigenvector, so \(Qq\geq0\). The effective real divisor
\(F=\sum q_iC_i\) is nef: its tested intersections on its support are
nonnegative, and those with any other curve are nonnegative. Adding a
sufficiently small positive multiple of \(F\) to \(P_D\) stays inside
(3.2), since every corresponding coefficient has positive room \(a_i\).
It improves the maximum, a contradiction. Therefore \(Q\) is negative
definite. \(\square\)

\textbf{Lemma 3.2 (inverse sign and numerical minimality).} For a
negative definite symmetric \(Q\) with nonnegative off-diagonal entries,
\(Q^{-1}\) has nonpositive entries. If a class \(\ell\) has a
decomposition \(\ell=p+[N]\) as in (3.1), with \(p\) nef and \(Q\)
negative definite, then in every numerical decomposition \[
 \ell=q_0+[M],\qquad q_0\text{ nef},\quad M\geq0,
                                                               \tag{3.3}
\] one has the actual divisor inequality \(M\geq N\).

\textbf{Proof.} Put \(A=-Q\). It is positive definite with nonpositive
off-diagonal entries. If \(Ax=b\geq0\) and \(x=x_+-x_-\) has disjoint
nonnegative parts with \(x_-\ne0\), then \[
 0\leq b\cdot x_-=x_+^TAx_- -x_-^TAx_-<0.
\] The first term is nonpositive and the second is strictly positive.
Hence \(x\geq0\). Apply this to the coordinate vectors \(b\) to obtain
\(A^{-1}\geq0\) and \(Q^{-1}\leq0\).

Write \(N=\sum a_iC_i\), and split \(M=\sum m_iC_i+M'\) with \(M'\)
having no component in the negative support. Orthogonality and nefness
of \(q_0\) give \[
 Q(a-m)_i-M'\cdot C_i\geq0.
\] Every \(M'\cdot C_i\geq0\). Therefore \(Q(a-m)\geq0\), and the
inverse sign implies \(a-m\leq0\). Thus \(M\geq N\). \(\square\)

Any two decompositions with these properties have the same negative
divisor, by applying minimality in both directions. In particular the
divisor \(N\) depends only on the numerical class, not on the chosen
effective representative, section, or ample perturbation. The positive
part of an \textbf{actual} bundle \(L\) is the actual rational bundle
\(P=L-\mathcal O(N)\), abbreviated to \(L-N\); it is not an unspecified
numerical representative.

\textbf{Lemma 3.3 (exact section division).} For a rational bundle
\(L=P+N\) with these properties, multiplication by the canonical section
is an actual isomorphism \[
 H^0(X,mP)\xrightarrow{\;\cdot s_{mN}\;}H^0(X,mL)
                   \quad(mP,mN,mL\text{ integral}).              \tag{3.4}
\] More generally every section of integral \(mL\) vanishes along
\(\lceil mN\rceil\).

\textbf{Proof.} The effective zero divisor \(T\) of a nonzero section of
\(mL\) is numerically \(mL\). Write \(T=\sum t_iC_i+T'\) outside the
negative support. Intersecting gives \[
 Q(t-ma)_i=-T'\cdot C_i\leq0.
\] Multiplying by \(Q^{-1}\leq0\) yields \(t\geq ma\). Since \(t_i\) are
integers, \(t_i\geq\lceil ma_i\rceil\). Dividing a section by the
indicated canonical section is globally a section of the actual residual
line bundle: its divisor is effective, and the division agrees in every
local frame. Conversely multiplication is injective and every residual
section gives an \(mL\) section. This proves (3.4), including the
zero-dimensional section spaces. \(\square\)

There is no finite-generation assumption on the section ring. The entire
base ideal of any chosen system is divided, not merely its divisorial
values.

\section{4. Pseudoeffective rational surfaces: finite support and the
perturbation
limit}\label{pseudoeffective-rational-surfaces-finite-support-and-the-perturbation-limit}

\textbf{Theorem 4.1 (canonical decomposition on the whole rational
pseudoeffective domain).} Every actual pseudoeffective rational \(L\)
has a unique decomposition \[
 L=P+N,
 \quad P\text{ actual nef rational},\quad N\geq0\text{ rational},
 \quad P\cdot C=0\ (C\subset\operatorname{Supp}N),
                                                               \tag{4.1}
\] with negative definite support when \(N\ne0\). If \(N_\epsilon\) and
\(P_\epsilon\) are the parts of \(L+\epsilon H\), for positive rational
\(\epsilon\), then \[
 0\leq N_\epsilon\leq N,\quad
 N_\epsilon\uparrow N\text{ coefficientwise as }\epsilon\downarrow0,
 \qquad c_1(P_\epsilon)\to c_1(P).                              \tag{4.2}
\] All negative supports in this family lie in one finite set of prime
curves. No effectivity or section of \(L\) is required.

\textbf{Proof.} For every \(\epsilon>0\), \(L+\epsilon H\) is an actual
big rational bundle. Indeed choose effective rational divisors \(D_j\)
with classes tending to \(c_1(L)\). On a tail \(L+\epsilon H-D_j\) is
actual ample. A divisible section of \(D_j\) multiplies the sections of
this ample difference into those of the actual \(L+\epsilon H\). Ample
Serre vanishing and Riemann--Roch give positive quadratic section
growth. Thus choose an effective representative of \(L+\epsilon H\) and
apply Lemmas 3.1--3.2.

If \(0<\epsilon'<\epsilon\), then \[
 L+\epsilon H=(P_{\epsilon'}+(\epsilon-\epsilon')H)+N_{\epsilon'}
\] is a nef-plus-effective decomposition. Minimality gives
\(N_\epsilon\leq N_{\epsilon'}\). Their supports are nested as
\(\epsilon\) decreases.

Let \(\rho=\dim N^1(X)_{\mathbb R}\). Each negative support has at most
\(\rho\) curves: a numerical dependence between these curves would give
a kernel vector of their negative definite intersection matrix. The
nested union also has at most \(\rho\) curves. Otherwise choose
\(\rho+1\) distinct curves in it and one rational \(\epsilon\) smaller
than the finitely many parameters exhibiting them; its negative support
would contain all of them. Hence the union is finite. Choose
\(\epsilon_0\) for which it is the entire support, and retain its
negative definite matrix \(Q\).

For \(0<\epsilon\leq\min(1,\epsilon_0)\), \[
 0\leq N_\epsilon\cdot H
       \leq (L+\epsilon H)\cdot H\leq L\cdot H+H^2.
\] Every supported prime curve has positive intersection with \(H\).
Therefore each nondecreasing coefficient of \(N_\epsilon\) is bounded.
Let its limit define \(N\). Its supported coefficients are positive
because each is already positive at \(\epsilon_0\). The closed nef cone
gives \(P=L-N\) nef. On the entire finite support,
\(P_\epsilon\cdot C_i=0\) for \(\epsilon\leq\epsilon_0\), so
\(P\cdot C_i=0\). In vector notation \[
 Q a=(L\cdot C_i)_i.
\] \(Q\) is invertible with rational entries and the right side is
rational. Thus \(a\) is rational, and \(P\) is an actual rational
bundle. Its negative matrix is \(Q\). Lemma 3.2 applies and proves
uniqueness and numerical minimality. It also gives \(N_\epsilon\leq N\)
by applying it to \(L+\epsilon H=(P+\epsilon H)+N\). These are exactly
(4.2). The empty-support case gives \(N=0\) without a matrix.
\(\square\)

For a big \(L\), one can alternatively start with Lemma 3.1 on an
effective representative. Minimality immediately gives
\(N_\epsilon\leq N\), so the fixed-support perturbation argument of
root183 is valid. The proof above explains why the absence of a section
of a pseudoeffective \(L\) does not create an infinite-support limit on
a surface. Its finite numerical rank and negative definiteness are
essential.

\section{5. The entire original-carrier surface
cone}\label{the-entire-original-carrier-surface-cone}

For a rational effective divisor \(N\), define its canonical potential
\[
 \sigma_N=\frac1q\log|s_{qN}|,
                   \quad qN\text{ integral}.                    \tag{5.1}
\] Tensor compatibility of the canonical section makes this independent
of \(q\). It is an actual one-section FS potential on \(N\), has generic
value zero, and is finite at every divisorial valuation. Its poles are
the analytic points whose support lies in \(\operatorname{Supp}N\).

\textbf{Theorem 5.1 (entire pseudoeffective rational surface cone).} For
every actual pseudoeffective rational \(L\) and its canonical
decomposition of Theorem 4.1, \[
 \operatorname{PSH}(c_1(L))
       =\{\sigma_N+v:v\in\operatorname{PSH}(c_1(P))\}.           \tag{5.2}
\] The equality is literal at every analytic point. Every function on
the right has one decreasing actual FS net, with numerical classes
tending to \(c_1(L)\), exact negative-part factor, and real weights. The
translation induces a homeomorphism of the entire normalized cones in
their original divisorial topologies. Those cones are compact. The
entire \(L\)-cone is decreasing closed for arbitrary nets with finite
limiting generic value.

\textbf{Proof: construction from the right.} Lemma 2.3 gives one finite
continuous actual FS net \(p_a\searrow v\) on \(P+\delta_aH\), with
\(\delta_a\to0\). Tensor with the canonical section of \(N\), clearing
degrees at each stage by all products. This gives \[
 \sigma_N+p_a\in\operatorname{FS}^{\rm gf}(L+\delta_aH).         \tag{5.3}
\] Its local base ideal is precisely the corresponding integral power of
the negative-part ideal: the second section system generates. Its
normalized vanishing orders are those of \(N\), not a small coefficient
of that divisor. At a pole of \(\sigma_N\) every stage equals
\(-\infty\). Elsewhere adding its finite value commutes with the
decreasing limit. The generic limit is finite. Formula (5.3) proves
genuine one-net membership.

\textbf{Proof: every original defining net can be divided.} Let
\(u\in\operatorname{PSH}(c_1(L))\), normalized, and fix rational
\(\epsilon>0\). By Lemma 2.2 its original FS stages, on a tail and with
unchanged functions, can be tensored into the \textbf{same actual}
\(L_\epsilon=L+\epsilon H\). Before dividing, raise each finite system
by all products to a degree also clearing \(N_\epsilon\). Lemma 3.3
divides every section by \(s_{mN_\epsilon}\), leaving a finite actual
system \(g_i\in\operatorname{FS}^{\rm gf}(P_\epsilon)\) with the
original real weights. At every analytic point, \[
 f_i=\sigma_{N_\epsilon}+g_i.                                  \tag{5.4}
\] This is an identity of products evaluated before logarithms; at a
shared zero both sides are \(-\infty\). It does not define \(g_i\) by
subtracting two infinities.

At every divisorial valuation, \(g_i=f_i-\sigma_{N_\epsilon}\), so these
values decrease. Choose actual ample \(A\) with \(A-P_\epsilon\) ample
and tensor all \(g_i\) unchanged into \(A\). Given two comparable
indices, their maximum is an actual FS function on \(A\): clear degrees
by all section powers and take the union of the systems with their
weights. It equals the earlier \(g_i\) on divisorial coordinates. Both
functions are already genuine in the same actual ample cone, so \(G_0\)
all-point determination proves their equality everywhere. Hence the
original \(g_i\) decrease at all points. Their limiting generic value is
\(u(x_0)=0\).

Proposition 2.4 yields an actual
\(v_\epsilon\in\operatorname{PSH}(c_1(P_\epsilon))\) with one defining
net and \[
 v_\epsilon(z)=u(z)-\sigma_{N_\epsilon}(z)
                     \quad(z\text{ divisorial}).               \tag{5.5}
\] It is normalized. The construction already proved from the right
makes \(\sigma_{N_\epsilon}+v_\epsilon\) genuine in \(c_1(L_\epsilon)\).
The unchanged original net makes \(u\) genuine there too. Embed both
into one dominating actual ample cone. Equation (5.5) and determination
identify them \textbf{at all points}, including the poles.

\textbf{Proof: the varying nef limit.} Let rational
\(\epsilon\downarrow0\). By (4.2), the classes \(c_1(P_\epsilon)\) tend
to \(c_1(P)\). Lemma 2.5 gives a convergent subnet of the normalized
actual \(v_\epsilon\) and a genuine normalized
\(v\in\operatorname{PSH}(c_1(P))\). At each divisorial \(z\), (4.2) and
(5.5) give the prescribed finite limit \[
 v(z)=u(z)-\sigma_N(z).                                         \tag{5.6}
\] Here \(\sigma_{N_\epsilon}(z)\) is a finite linear combination of the
finitely many component evaluations, so its coefficient limit is
legitimate. We have not interchanged two all-point limits. The
constructed \(\sigma_N+v\) is genuine in \(c_1(L)\) by (5.3). Both it
and the already genuine \(u\) embed into one actual ample cone. Equation
(5.6) therefore identifies them at every point. This proves (5.2).

\textbf{Proof: compactness and decreasing closure.} On divisorial
coordinates the translation is addition by the fixed finite coordinates
of \(\sigma_N\). It is injective: identical images give identical nef
coordinates, and a common ample embedding identifies the two genuine nef
functions at all points. Its inverse is coordinate translation as well,
and generic normalization is preserved. Proposition 2.4 therefore gives
compactness of the entire normalized \(L\)-cone.

For a decreasing net \(u_i=\sigma_N+v_i\) with finite limiting generic
value, the \(v_i\) decrease on divisorial coordinates. Embed any
comparable pair into a common ample cone. Its genuine maximum is
obtained by maxima of its two finite continuous ample regularizing nets;
it has a decreasing same-bundle defining net. Determination promotes the
coordinate order to all-point order. Thus \(v_i\) is an all-point
decreasing nef net. Proposition 2.4 supplies \(v=\inf_i v_i\) and one
actual defining net. At the poles of \(\sigma_N\), every \(u_i\) and
\(\sigma_N+v\) equals \(-\infty\); elsewhere addition commutes with the
infimum. Consequently \(\inf_i u_i=\sigma_N+v\), and (5.3) supplies the
required \textbf{one} actual defining net in the original \(L\) class.
\(\square\)

This verifies the entire-cone assertion of root183, with no additional
hypotheses beyond the retained smooth surface, field, rationality, and
ample analytic grant. It also proves the stated pseudoeffective
extension. It would be incorrect to say that the entire big surface cone
has globally finite defining FS stages: its stages in (5.3) keep \(N\)
exactly.

\textbf{Corollary 5.2 (intrinsic minimal reference and actual floors).}
In the surface domain, \[
 V_L:=\sup\{u\in\operatorname{PSH}(c_1(L)):u\leq0\}=\sigma_N.
                                                               \tag{5.7}
\] The supremum is attained and usc. If \(u=\sigma_N+v\), then \[
 \max(u,V_L-j)=\sigma_N+\max(v,-j)\quad(j\geq0).                 \tag{5.8}
\] Both sides are actual genuine functions; this identity is valid at
shared poles.

\textbf{Proof.} For \(u\leq0\), \(v(x_0)=u(x_0)\leq0\), and (1.3) gives
\(v\leq0\). Hence \(u\leq\sigma_N\). The choice \(v=0\), genuine by
Lemma 2.3, gives \(\sigma_N\) itself. The nef maximum in (5.8) is
genuine by a common ample regularization and Lemma 2.3. Literal extended
addition proves (5.8), including where \(\sigma_N=-\infty\). \(\square\)

No assertion that all finite-energy potentials have bounded difference
from \(V_L\) is implicit in (5.8). The entire positive-volume
finite-energy space has not yet been constructed here.

\section{6. Varying-class surface closure without an extended-real
flattening
lemma}\label{varying-class-surface-closure-without-an-extended-real-flattening-lemma}

\textbf{Proposition 6.1.} For an actual pseudoeffective rational surface
bundle \(L\), \[
 \operatorname{PSH}(c_1(L))
   =\bigcap_{\epsilon\in\mathbb Q_{>0}}
                      \operatorname{PSH}(c_1(L+\epsilon H)).      \tag{6.1}
\] If actual rational classes \(\theta_i\to c_1(L)\) and genuine
\(u_i\in\operatorname{PSH}(\theta_i)\) decrease at all points with
finite limiting generic value, then their all-point limit is genuine in
\(c_1(L)\) and has one defining net. This is the varying-class closure
property at every rational pseudoeffective surface endpoint.

\textbf{Proof.} Forward inclusion of (6.1) is Lemma 2.2. Conversely let
one actual function \(u\) belong to every cone on the right, and
normalize it. Theorem 5.1, applied to each \(L+\epsilon H\), provides
\(u=\sigma_{N_\epsilon}+v_\epsilon\) with genuine normalized nef
\(v_\epsilon\). Lemma 2.5 and (4.2) give a genuine normalized nef limit
\(v\) on a subnet, whose divisorial coordinates are
\(u(z)-\sigma_N(z)\). The constructed \(\sigma_N+v\) is genuine in the
\(L\) cone. The function \(u\) is already genuine in, for example, the
\(L+H\) cone; both embed into one actual ample cone. Determination makes
their equality all-pointwise. Formula (5.3) witnesses membership with
one net.

Now fix \(\epsilon>0\). For all sufficiently large \(i\),
\(c_1(L+\epsilon H)-\theta_i\) is ample. Lemma 2.2 puts the same
\(u_i\), on that tail, into the entire genuine \(L+\epsilon H\) cone.
Its decreasing closure from Theorem 5.1 puts the same all-point limit
there. This holds for each rational \(\epsilon\), so (6.1) applies. Its
proof and (5.3) give the single actual FS defining net. \(\square\)

The cone compactness in Theorem 5.1 was obtained as an image of an
independently proved nef compact cone. It did not assume Proposition 6.1
or general big compactness. There is no circular appeal to the disputed
BJ Lemma 4.6 or its dependent nonample compactness theorem.

\section{7. Algebraic volume and raw mass: what is and is not
constructed}\label{algebraic-volume-and-raw-mass-what-is-and-is-not-constructed}

\textbf{Theorem 7.1.} For every actual pseudoeffective rational surface
bundle \(L=P+N\) of Theorem 4.1, \[
 \operatorname{vol}_X(L)=\operatorname{vol}_X(P)=P^2.             \tag{7.1}
\] In the big case \(P^2>0\). Thus any full-mass nonpluripolar equation
in this domain must use the raw mass \(P^2\), which need not equal
\(L^2\).

\textbf{Proof.} First prove an upper bound that also handles bounded
twists and residue degrees. For a fixed actual integral twist \(B\) and
sufficiently divisible \(m\) for \(P\), choose \(tH-B\) with a nonzero
section. Then \[
 h^0(mP+B)\leq h^0(mP+tH).
\] For rational \(\delta>0\), multiply further by a nonzero section of
\(\lceil m\delta\rceil H\) for all large \(m\). The resulting bundle is
\[
 mP+(t+\lceil m\delta\rceil)H.
\] After separating the finitely many residue classes modulo a
denominator for \(P\) and \(\delta\), it is a growing integral multiple
of the fixed ample rational \(P+\delta H\) plus one of finitely many
fixed actual twists. Serre vanishing and Riemann--Roch give \[
 \limsup_m \frac{2h^0(mP+B)}{m^2}\leq(P+\delta H)^2.
                                                               \tag{7.2}
\] The same argument for any permissible arithmetic progression is
valid. Let \(\delta\downarrow0\) to get upper bound \(P^2\). Also
\(P^2\geq0\), since \(P+\delta H\) is ample.

If \(P^2>0\), then \(P\cdot H>0\). Otherwise choose small rational
\(a>0\) with \(H-aP\) ample; nefness would give
\(0\leq P\cdot(H-aP)=-aP^2\), a contradiction. For large divisible
\(m\), \(K_X-mP\) has negative intersection with \(H\), so it has no
nonzero section. Serre duality and Riemann--Roch imply \[
 h^0(mP)\geq\chi(mP)=\tfrac12m^2P^2+O(m).
\] Together with (7.2) this gives volume \(P^2\). If \(P^2=0\), the
upper bound and nonnegative section counts give volume zero.

Lemma 3.3 gives section-space equality for all degrees divisible by both
\(L\) and \(N\). To verify the upper bound for every degree clearing
\(L\), and not merely one subsequence, every section divides by
\(\lceil mN\rceil\). Choose \(q\) clearing \(L,N,P\) and write the
permissible \(m=aq+r\). Then \[
 mL-\lceil mN\rceil=aqP+B_r,
 \qquad B_r=rL-\lceil rN\rceil,
\] with only finitely many actual integral residual twists \(B_r\).
Apply (7.2) with these twists and with growing degree \(aq\);
\(aq/m\to1\). This bounds the entire section-volume limsup of \(L\) by
\(P^2\). The divisible section-space equality and the lower bound for
\(P\) give the reverse inequality when \(P^2>0\); when \(P^2=0\)
nonnegativity suffices. This proves (7.1) without an unproved
arithmetic-subsequence interchange. \(\square\)

For example, on the Hirzebruch surface with section \(S^2=-2\), fiber
\(F^2=0\), and \(S\cdot F=1\), the actual divisor \(L=2S+3F\) has \[
 N=\tfrac12S,\quad P=\tfrac32S+3F,\quad L^2=4,
                     \quad\operatorname{vol}(L)=P^2=\tfrac92.
\] For \(H=S+3F\) and \(0<\epsilon\leq1\),
\(N_\epsilon=(1-\epsilon)S/2\). This is an explicit section-division and
perturbation diagnostic, not a finite proof of the general theorem.

The canonical trivial nef metric on \(P\) has the ordinary product-model
raw top product \(P^2\delta_{x_0}\). Removing the fixed canonical \(N\)
factor from the reference metric leaves this actual moving model metric.
This verifies the reference raw moving mass. Theorem 7.1 alone does
\textbf{not} construct a measure \(\operatorname{MA}_L(u)\) for every
unbounded relative finite-energy \(u\). Assigning the desired mass to an
inverse-selected measure would not supply that operator.

\section{8. Two explicit zero-volume equation
ranges}\label{two-explicit-zero-volume-equation-ranges}

\subsection{8.1. Entire numerical-zero class in every positive
dimension}\label{entire-numerical-zero-class-in-every-positive-dimension}

\textbf{Theorem 8.1.} For an actual numerically trivial rational \(B\)
on the original \(X\) of any dimension \(n\geq1\), \[
 \operatorname{PSH}(c_1(B))=\mathbb R.
                                                               \tag{8.1}
\] With canonical reference zero, the actual entire finite-energy space
is \(\mathbb R\), raw energy is zero, and the genuine raw model operator
is zero. Its only positive target is the zero Radon measure. All
solutions are unique modulo constants.

\textbf{Proof.} Lemma 2.3 puts constants in the cone, even if \(B\) is
nontorsion and has no sections. Let \(u\) be normalized in the cone.
Integer tensor powers of its original defining net show directly that
\(mu\) is genuine in the zero numerical class for each positive integer
\(m\). Lemma 2.2 embeds these unchanged normalized functions into one
actual ample \(H\) cone. At a fixed divisorial \(z\), the continuous
finite evaluation on its compact normalized cone has a finite lower
bound \(-C_z\). Thus \[
 -C_z\leq m u(z)\leq0 \qquad(m\geq1).
\] It follows that \(u(z)=0\). Ample determination for the already
genuine \(u\) and zero in the same ample cone gives \(u=0\) at every
point. This proves (8.1).

Every section of a divisible numerically trivial bundle is either zero
or nowhere vanishing. Indeed a nonzero effective zero divisor would have
strictly positive intersection with \(H^{n-1}\) but numerical class
zero. A nonzero section therefore trivializes that power, and proper
integrality gives dimension at most one for its section space. Its
volume is zero.

On the product model, the actual canonical metric has raw top product
\((B^n)\delta_{x_0}=0\). Constants do not change it. The raw energy
normalized at zero obeys \(E(c)=c(B^n)=0\); the model intersection gives
this identity directly. Since all genuine functions are constants, its
entire finite-energy space is exactly \(\mathbb R\), with no unproved
extension to other functions. The ample products on \(B+\epsilon H\) at
a constant also converge to zero, with masses \(\epsilon^n(H^n)\).

A positive Radon measure of raw mass zero is the zero measure. The
intrinsic dual energy is \(\sup_c(E(c)-\int c\,d0)=0\). Comparison is
the equality of two zero integrals. Every constant solves the genuine
zero equation, and any two solutions differ by a constant. The
normalized cone is a compact single point. For every continuous obstacle
\(f\), its envelope is the constant \(\min f\) and its outer energy is
zero; both energy derivatives are zero. No probability normalization or
division by volume occurs. \(\square\)

\subsection{8.2. Surface classes with numerically zero positive
part}\label{surface-classes-with-numerically-zero-positive-part}

\textbf{Theorem 8.2 (flat positive-part relative equation).} For the
surface domain of Theorem 5.1, assume additionally \(c_1(P)=0\) in the
canonical decomposition \(L=P+N\). Then \[
 \operatorname{PSH}(c_1(L))=\{\sigma_N+c:c\in\mathbb R\}.
                                                               \tag{8.2}
\] With intrinsic reference \(V_L=\sigma_N\), its entire relative
finite-energy space is this set, its genuine raw nonpluripolar top
product is zero, and its relative energy is zero. The sole positive
target is zero; the equation is unique modulo constants. This includes
nonzero pseudoeffective classes, for example the exceptional divisor on
a blown-up surface.

\textbf{Proof.} Equation (8.2) is Theorem 5.1 and Theorem 8.1. It
remains to specify the relative product and energy rather than silently
using a signed fixed-divisor self-intersection.

Let \(U=X\setminus\operatorname{Supp}N\). In the actual local frame
obtained by the nonvanishing canonical \(N\) section on \(U\), the
metric with scalar potential \(\sigma_N+c\) is precisely the canonical
metric of the actual \(P\), shifted by \(c\). The \(N\) factor has been
removed as an actual section, not by scalar subtraction at a pole. The
product-model curvature product of this numerically trivial metric is
zero on \(U^{\rm an}\). Define its nonpluripolar extension to the
original \(X^{\rm an}\) by giving zero mass to the reference pole set
\((\operatorname{Supp}N)^{\rm an}\). The result is the zero measure.
This is the extension of the actual local intersection product,
independent of any target and of any choice of lift. On the nef positive
part, it is also the weak limit of actual ample constant products on
\(P+\epsilon H\), whose masses are \(\epsilon^2H^2\to0\).

The relative energy difference of two translated metrics is the usual
integral of their constant difference against the sum of the local
positive curvature products, extended with zero polar mass. Every term
is zero. Normalize \(E_L(\sigma_N)=0\); thus \(E_L(\sigma_N+c)=0\). This
constructs the actual relative energy on the \textbf{entire} cone (8.2).
The relative floors in (5.8) stay in this same cone and give no further
endpoints. Its entire finite-energy space is therefore (8.2).

At the reference poles, integrals against the only target zero are
interpreted on \(U^{\rm an}\), with no \(\infty-\infty\) convention.
Intrinsic target mass is \(\operatorname{vol}(L)=P^2=0\) and relative
dual energy is zero. Comparison integrals vanish. Every translated
reference is an actual solution and any two differ by a constant.
Normalization gives one actual point, so attainment is literal. For a
continuous perturbation \(f\), the largest cone member below
\(\sigma_N+c+tf\) is \(\sigma_N+c+\inf_{U^{\rm an}}tf\) (with equality
at reference poles automatic). Its outer relative energy is still zero;
both derivative signs agree with the zero product. \(\square\)

This explicit construction is confined to (8.2). It does not define an
operator for the entire positive-volume cone by postulating a transport
from an unconstructed nef finite-energy theory. At zero mass, the usual
unnormalized domination statement fails even for constants: \(c<d\) but
a zero measure sees no set distinguishing them. Uniqueness here comes
from cone rigidity modulo constants, not from such domination.

\section{9. A nonzero zero-volume class where model uniqueness
fails}\label{a-nonzero-zero-volume-class-where-model-uniqueness-fails}

Let \(X=\mathbb P^1\times\mathbb P^1\), write \(A=\mathcal O(1,0)\),
\(B=\mathcal O(0,1)\), and take \(L=A\), \(H=A+B\). Then \(L\) is actual
nef, \(N=0\), \(L^2=0\), and the minimal reference is zero. Let \(s,t\)
generate the first \(\mathcal O(1)\) and set, on the original \(X\), \[
 f=\max(\log|s|,\log|t|-1).
                                                               \tag{9.1}
\] It is a finite continuous actual FS potential on \(A\). Both \(0\)
and \(f\) are normalized genuine potentials. Along
\(a\operatorname{ord}_{s=0}\), \(f=-\min(a,1)\); on the opposite fiber
valuation it is zero. Thus \(f\) is nonconstant.

The model producing (9.1) is the pullback of the curve FS model. On its
product with the second \(\mathbb P^1\), every curvature line bundle for
this metric is pulled from the two-dimensional curve test configuration.
Its square restricted to a central component of the surface model is
zero, and its three-fold energy intersection is zero. Hence the ordinary
genuine nef \textbf{model} top products of \(0\) and \(f\) are zero and
their raw model energies relative to zero are zero. This is a dimension
calculation of intersection products, not an inverse measure
construction.

The same conclusion follows for the raw positive top-product limit from
the task-granted ample theory: Lemma 2.1 puts both functions unchanged
in the actual ample \(A+\epsilon H\) cone; each of its positive measures
has mass \[
 (A+\epsilon H)^2=2\epsilon+2\epsilon^2\longrightarrow0.
\] Therefore their integrals against any continuous test tend to zero,
bounded by its sup norm times this mass. Any zero-volume equation theory
retaining these bounded model metrics, their actual finite model
energies, and their genuine raw top products has two normalized
solutions of the zero target. Uniqueness modulo constants consequently
fails in this class. A theorem on all nonzero zero-volume
pseudoeffective classes cannot merely repeat the positive-volume
uniqueness statement.

No full unbounded operator on this class is asserted. The bounded model
counterexample already proves the obstruction to the proposed
unrestricted zero-mass uniqueness requirement.

\section{Appendix B. Entire intrinsic nef operator, weighted limits,
outer derivatives and
dual}\label{appendix-b.-entire-intrinsic-nef-operator-weighted-limits-outer-derivatives-and-dual}

The following complete mathematical derivation is retained as supporting
proof, with its earlier scope statements. It is not a historical-verdict
grant. Equation numbers in this appendix are local to it. The main
candidate and the fresh review criteria govern the new combined scope.

\section{Intrinsic finite-energy Monge--Ampere products in nef
classes}\label{intrinsic-finite-energy-mongeampere-products-in-nef-classes}

This paper states and proves the exact nef operator and
variational-support theorem independently accepted in Xean review124.
The field is algebraically closed, trivially valued, of characteristic
zero and arbitrary cardinality; the original variety is smooth integral
and projective of dimension \(n\geq1\). The datum is any actual nef
rational line bundle \(L\), with no semiampleness or finite-generation
assumption. All potentials and measures remain on the entire original
\(X^{\mathrm{an}}\).

The reference is the canonical zero metric. The raw mass is
\(V=c_1(L)^n\). On the entire genuine numerical cone define

\[E_\theta(u)=\inf_{k\geq1}E_\theta(\max(u,-k)),\qquad
\mathcal E^1(\theta)=\{u\in\operatorname{PSH}(\theta):E_\theta(u)>-\infty\}.\]

The theorem constructs a target-independent positive Radon top operator
and common-class mixed products of mass \(V\) on this whole intrinsic
domain, including endpoints outside the auxiliary ample finite-energy
vector space. It proves reference-floor total-variation convergence,
integrability of every intrinsic test, weighted arbitrary decreasing-net
convergence against all original continuous tests, both signs of the
original continuous-direction outer derivative, weak upper
semicontinuity and compact energy-cost sublevels. At \(V>0\), the
intrinsic finite-dual condition is equivalent to integration of every
genuine finite-energy test, every constructed output satisfies it, and
its variational objective is coercive.

At \(V=0\), the positive top and common-class mixed operators, energies
and weighted products vanish. No probability normalization, inverse
equation or uniqueness is claimed in this range. This paper proves no
all-target attainment or uniqueness theorem at positive volume either.
Thus its reviewed scope is a supporting theorem for original P16.5; it
does not close that original problem, and supplies no general nonnef
big-class extension.

The only supplied equation input is the exact actual-ample G0 theorem,
used solely at positive actual ample perturbations. The published
model-pairing, envelope and numerical inputs are named in each proof
part and remain in their stated domains. They are imported mathematical
foundations, not claims of recursively proving every cited theorem. The
four proof parts below preserve the frozen mathematical bodies
independently reconstructed in review124. Their earlier draft-status
sentences and chronological remaining-obligation paragraphs are retained
as provenance; the present theorem and Part V specify the accepted final
scope.

\subsection{0. The ample input and the finite-continuous closure
replacement}\label{the-ample-input-and-the-finite-continuous-closure-replacement}

The supplied ample theorem concerns an actual ample rational bundle on
the retained original characteristic-zero projective variety. It
constructs its genuine finite-energy metric operator, raw energy and
continuous-obstacle envelopes, supplies compactness and all-point
divisorial determination, and proves both continuous-direction outer
energy derivatives. It gives no assertion for a nonample class. Its
exact full mathematical source is the previously accepted P16.1 paper;
the independently frozen review packet includes that entire body. Within
the present proof every invocation has an actual ample class, such as
\(\theta+\epsilon h\) with \(\epsilon>0\).

For clarity, the closure step avoids an arbitrary extended-valued
flattening lemma. Genuine nef membership equals the intersection of the
genuine cones for \(\theta+\epsilon h\). The forward section embeddings
preserve the scalar, since generating canonical section norms have
maximum zero. For the reverse direction choose finite continuous ample
FS majorants of a potential in every positive perturbed cone, allowing
positive constant shifts. They form a single directed decreasing family:
for any finite collection of old continuous majorants, their finite
continuous minimum is an upper bound; with a positive strict tolerance,
eventual approximation below that minimum holds on an open cover, and
compactness provides one common later approximant. Refine the positive
perturbation at the same step. Infimizing tolerances and majorants gives
the desired original genuine all-point limit, including its poles. This
uses finite continuous old minima, never a general varying-class
flattening principle for arbitrary discontinuous functions. Closed
intersections in one normalized ample compact cone provide the genuine
normalized nef compactness and the envelope passage used below.

\subsection{I. Bounded genuine metric products and continuous
obstacles}\label{i.-bounded-genuine-metric-products-and-continuous-obstacles}

This is a new unreviewed own deduction for original P16.5. It is not
supplied to either active frozen native119 or native121. Original
closure and completion counts do not change. The deduction advances the
general nef case without semiampleness, while preserving the distinction
between bounded potentials and the larger unbounded nef finite-energy
class.

\subsubsection{1. Actual scope and the corrected source
boundary}\label{actual-scope-and-the-corrected-source-boundary}

Let k be algebraically closed, trivially valued and of characteristic
zero, with arbitrary cardinality. Let X be smooth, integral and
projective of dimension n\textgreater=1. Retain an actual nef rational
line bundle L, its numerical class theta, and its canonical trivial
reference zero. Fix an actual ample rational bundle H with class h. All
potentials below belong to the genuine numerical cone of BJ
Definition4.1; no quotient cone or inverse-lift operator is defined.

The primary source is Boucksom--Jonsson, arXiv:1801.08229v4, header
March2022, 143pages. Theorem5.20 itself states the envelope property for
\textbf{every nef theta}, not just ample theta, under the retained
smoothness and characteristic assumptions. Its printed proof reduces to
ample rational classes through Lemma5.9; that reduction does not narrow
the theorem's conclusion. Complete printed statement/proof units checked
here are local literal LF2980--2985 and3027--3062, and fresh native-PDF
lines4269--4274 and4321--4363. Lemma5.9's complete proof is already
checked at LF2801--2842. Published recursive inputs in the
multiplier-ideal argument are retained, not independently recursively
certified.

Two distinctions matter. Lemma5.17's equivalence with continuity of
continuous-obstacle envelopes is formulated for ample classes.
Theorem5.20 does not by that citation alone prove continuity of every
nef continuous-obstacle envelope. Remark7.16 further warns that the
floor-defined E1(theta) can strictly contain the ambient finite-energy
vector space intersected with PSH(theta). Theorems7.14/7.18 and
Proposition7.19 therefore cannot simply be applied to every member of
E1(theta).

The exact whole accepted ample G0 proof is
\nolinkurl{/workspace/research/root-checks/public-P16.1-round-92-final-reader/argument.md},
SHA256
\nolinkurl{6cd26d798b63a388bd9a5e6dcd0583a7cd9c7417648efe390f4d4d00d28646b9}.
Only its actual ample projective domain is a grant. Its raw energy,
genuine metric operator, energy concavity, continuous envelopes, and
both continuous-direction outer derivatives are used at theta+epsilon h.
No nonample equation or unbounded nonample continuity is granted.

\subsubsection{2. Bounded genuine potentials fit the actual published
product
domain}\label{bounded-genuine-potentials-fit-the-actual-published-product-domain}

Write B(theta)=PSH(theta) intersect bounded functions. This is a
statement about the entire genuine numerical bounded cone, not a newly
defined smaller test family. Since theta is nef, zero and every constant
are genuine theta-psh by BJ4.3. If u belongs to B(theta), then u belongs
to PSH(theta+h), since zero is h-psh and numerical addition is BJ4.7. It
is bounded, so its ample energy is finite by monotonicity between its
constant bounds. Thus u belongs to the actual ambient vector space
E\textasciitilde1 of Definition7.12.

Theorems7.14 and7.18 now apply with their stated hypotheses. Define the
raw bounded energy and mixed measures by their actual energy-pairing
formulas:

\[E_\theta(u)=\frac{(\theta,u)^{n+1}-(\theta,0)^{n+1}}{n+1},\qquad
M_\theta(u_1,\ldots,u_n)=\bigwedge_{i=1}^n(\theta+dd^cu_i).\]

The total mass is theta\^{}n, and M is positive for genuine bounded
theta-psh inputs, by Theorem7.18 and Proposition7.19(i). The energy
shifts by c theta\^{}n under u -\textgreater{} u+c.~At a finite PL/model
stage, Proposition7.19(ii) gives the actual sum of component Dirac
measures with coefficients

\[\operatorname{ord}_F(\mathcal X_0)\prod_i(\theta_{\mathcal X}+D_i)|_F.\]

These are actual model metric intersection products with the given
reference, component multiplicities and bundle class. On continuous
potentials the published CLD compatibility of Remark7.23 applies.
Individual current factors are not asserted to have been constructed on
all bounded inputs merely from the wedge notation.

One can also recover the same bounded theta-products as limits of actual
positive ample operators before taking any limit. For
epsilon\textgreater0 rational, each u\_i remains genuinely
(theta+epsilon h)-psh. The actual G0 ample operator on L+epsilon H is
defined, positive, and has raw mass (theta+epsilon
h)\^{}n.~Multilinearity of the published mixed measures makes its
pairing with each continuous test a finite polynomial in epsilon, with
the constant term M\_theta(u\_1,\ldots,u\_n). Thus the positive ample
products converge weakly to the displayed genuine theta-product.
Uniqueness in Theorem7.18 makes this independent of H. The limit is
identified with an already constructed original metric product; it is
not the definition of an arbitrary inverse lift.

For a decreasing net u\_i in B(theta) with bounded genuine limit u, any
tail has common upper and lower bounds. All terms lie in E1(theta+h),
and Proposition7.19(iv) yields the original bounded mixed weak
continuity. Theorem7.14 gives bounded-energy continuity for these nets.
Integrability and the energy-pairing identity for an ambient
E\textasciitilde1 test are Proposition7.19(v). This proves no weighted
theorem for a test outside that stated ambient domain.

\subsubsection{3. Genuine envelopes at the nef
boundary}\label{genuine-envelopes-at-the-nef-boundary}

For continuous f on the original compact X\^{}an, let P\_theta f be the
pointwise supremum of genuine theta-psh functions bounded by f.~Constant
competitors show that it lies between min f and max f.~The nef envelope
property makes the upper regularization theta-psh; continuity of f makes
that regularization \textless=f, so it is itself a competitor. Hence
P\_theta f equals that regularization, is genuinely theta-psh, and is
bounded. This argument does not assert its continuity.

Put theta\_epsilon=theta+epsilon h and q\_epsilon=P\_theta\_epsilon
f.~The actual ample theory makes q\_epsilon continuous. As epsilon
decreases to zero, q\_epsilon decreases pointwise. The common constant
bounds make the limit generically finite, and BJ(PSH2) makes its limit
theta-psh. That limit is \textless=f.~Conversely every theta-psh
competitor is theta\_epsilon-psh, so P\_theta f\textless=q\_epsilon.
Therefore

\[q_\epsilon\downarrow q_0=P_\theta f.\]

All q\_epsilon, for epsilon\textless=1, lie in the one actual ample cone
PSH(theta+h) and have the same two constant bounds. Theorems7.14/7.18
and Proposition7.19 therefore apply to this decreasing family within
their proper domain. Polynomial expansion, coefficient by coefficient,
proves

\[E_{\theta_\epsilon}(q_\epsilon)\longrightarrow E_\theta(q_0),\qquad
M_{\theta_\epsilon}(q_\epsilon)\longrightarrow M_\theta(q_0).\]

Some individual theta-coefficients at positive epsilon may be signed;
positivity of those coefficients is neither used nor asserted. Their
continuous-test pairings converge by Proposition7.19(iv) in the common
ample finite-energy cone. The full operators on the left and the final
operator on the right are positive. Every finite polynomial coefficient
has a finite limit, so the positive powers of epsilon vanish without an
unproved uniform estimate over arbitrary E1(theta).

\subsubsection{4. Uniform continuity needed for
variation}\label{uniform-continuity-needed-for-variation}

Pointwise order and constants give the original envelope estimate

\[\|P_\theta f-P_\theta g\|_\infty\leq\|f-g\|_\infty.\]

For uniformly converging bounded genuine theta-psh potentials u\_j
-\textgreater{} u, their theta mixed measures and energy converge. Here
is a reduction to the published decreasing-net theorem that avoids any
unproved nonample strong topology. From any sequence converging in sup
norm, pass to a subsequence with errors \textless=4\^{}\{-j\}; adding
3\emph{4\^{}\{-j\} makes that subsequence decreasing pointwise to u.
Constants do not change curvature products. Energy changes only by
3}4\^{}\{-j\} theta\^{}n.~The decreasing convergence in the common ample
cone from Section2 therefore gives the desired convergence on that
subsequence. Every subsequence admits such a further subsequence, so the
entire sequence converges. For each fixed continuous measure test this
proves continuity on the sup-norm metric space; consequently it also
proves convergence for arbitrary sup-norm converging nets. It does not
apply a false monotone convergence theorem to arbitrary Borel nets.

\subsubsection{5. Both derivatives of continuous-obstacle
energy}\label{both-derivatives-of-continuous-obstacle-energy}

Define F\_theta(f)=E\_theta(P\_theta f) for f in C0(X\^{}an). It is
finite, concave and translation equivariant with raw coefficient
theta\^{}n.~Concavity follows from the genuine convex cone, envelope
order and bounded energy concavity; alternatively it follows from the
pointwise limit of the ample concave functionals in Section3.

We first give the needed ample supporting inequalities rather than
importing an unproved assertion about arbitrary obstacles. For an ample
class omega and continuous f, q=P\_omega f is continuous genuine. If
0\textless=s\textless=1, monotonicity gives P\_omega(q+s(f-q))=q. The
actual ample G0 outer derivative at q in the continuous direction f-q
therefore yields orthogonality:

\[\int(f-q)M_\omega(q)=0.\]

Energy concavity along genuine psh segments gives
E\_omega(r)\textless=E\_omega(q)+integral(r-q)M\_omega(q) for r=P\_omega
g. Since r\textless=g and M\_omega(q) is positive, the displayed
orthogonality proves the supporting inequality

\[F_\omega(g)-F_\omega(f)\leq\int(g-f)M_\omega(P_\omega f).\]

Use this inequality once at f and once at f+tg. For t\textgreater0 it
bounds the quotient between integral g M\_omega(P\_omega(f+tg)) and
integral g M\_omega(P\_omega f). Apply it with omega=theta\_epsilon and
pass epsilon down to zero using the actual energy and operator limits
from Section3. Thus

\[\int gM_\theta(P_\theta(f+tg))\leq
\frac{F_\theta(f+tg)-F_\theta(f)}t\leq
\int gM_\theta(P_\theta f).\]

The envelope Lipschitz estimate and Section4 make the left integral
converge to the right one as t down to zero. For t\textless0 the same
two supporting inequalities reverse under division; the two bounds have
the same limit as t increases to zero. Therefore both derivatives exist
and agree:

\[\left.\frac{d}{dt}\right|_{t=0}F_\theta(f+tg)=\int gM_\theta(P_\theta f),\qquad f,g\in C^0(X^{\rm an}).\]

In particular this proves both outer signs at every \textbf{continuous
genuine} theta-psh base potential u, by taking f=u. It proves the
continuous-obstacle formula for every nef actual L, without
semiampleness and including zero raw top mass. It does not infer a
derivative at an arbitrary unbounded member of the floor-defined
E1(theta), or even at an arbitrary bounded discontinuous base, from this
continuous-obstacle argument. There is no contradiction with
proposal131: that proposal uses a different lower-degree mixed energy of
positive mass, whereas the present top-degree energy has zero mass for
its product L and is constant.

\subsubsection{6. Remaining original
obligations}\label{remaining-original-obligations}

This supplies a genuine bounded top-degree operator, nef envelope
membership, and both derivatives for continuous obstacles beyond
semiample contractions. It does not supply the complete unbounded
E1(theta) operator with all original weighted/floor limits, intrinsic
target-space characterization, or variational attainment and uniqueness
for every finite-energy measure. The nef raw top mass is theta\^{}n;
proving all required section-volume and reference identifications in the
intended full general class remains part of the original problem.
General nonnef big references and general nonbig nonzero classes are
also still unresolved here.

The explicit published source domains are the important constraint.
Remark7.16 prohibits silently identifying the floor-defined E1(theta)
with E\textasciitilde1 intersect PSH(theta). A bounded approximation
theorem cannot be promoted to that larger class without an energy
estimate controlling the full limiting operator. All new conclusions
here require independent reconstruction. Original P16.5 remains open.

\subsubsection{7. Source/application
ledger}\label{sourceapplication-ledger}

\begin{itemize}
\tightlist
\item
  BJv4 LF2190--2237: genuine numerical definition, joint decreasing
  class closure, constants for nef theta. LF2330--2368: addition and
  uniform closure proof. Used for Sections2--4.
\item
  BJv4 LF2801--2842: Lemma5.9 and Theorem5.11 complete
  statements/proofs. LF2980--2985 and3027--3062: complete Theorem5.20
  statement and proof. Fresh native-PDF selected units are retained as
  source audit only. Used for nef envelope membership, not unproved nef
  continuity.
\item
  BJv4 LF2909--2935: complete Lemma5.17 statement/proof, whose ample
  domain is retained. Used only for theta\_epsilon continuous envelopes.
\item
  BJv4 LF3890--3955: actual ambient E\textasciitilde1 definition,
  complete Theorem7.14 and Lemma7.15 proofs, and Remark7.16.
  LF3980--4057: complete Theorem7.18 and Proposition7.19
  statements/proofs. Used only on verified bounded inputs in a common
  ample finite-energy cone.
\item
  Exact whole ample G0 body named in Section1: actual raw operator,
  energy concavity and both outer derivatives within its accepted ample
  domain. Every use at positive epsilon retains that domain. No original
  P16.5 or other unproved nonample premise is a grant.
\item
  Primary URL: https://arxiv.org/pdf/1801.08229v4 . Whole-book and
  recursive-foundation verification are not claimed. No finite numerical
  check certifies the deduction.
\end{itemize}

\subsection{II. The entire intrinsic floor operator and integrable
weights}\label{ii.-the-entire-intrinsic-floor-operator-and-integrable-weights}

This is a new unreviewed own deduction, not a grant or an original
completion. Neither active frozen native119 nor native121, nor prepared
frozen review122, receives this later material. The construction
addresses the larger floor-defined nef E1 class without assuming that
its members belong to an auxiliary ample finite-energy space. Arbitrary
decreasing-family continuity, a full variational inverse and general
original P16.5 remain open.

\subsubsection{1. Actual data and verified bounded input
domain}\label{actual-data-and-verified-bounded-input-domain}

Retain an algebraically closed trivially valued field k of
characteristic zero and arbitrary cardinality, a smooth integral
projective n-fold X with n\textgreater=1, an actual nef rational bundle
L with class theta, and its canonical reference zero. Put
V=theta\^{}n\textgreater=0. Use the genuine numerical PSH(theta) of BJv4
Definition4.1. Its constants and finite maxima are genuine. No
semiampleness is assumed.

For bounded u in PSH(theta), membership in PSH(theta+h) for an actual
ample H follows by adding the genuine zero h-potential. Boundedness
makes u finite-energy in that actual ample class, so u is in the
published ambient E\textasciitilde1. Theorems7.14/7.18 and
Proposition7.19 therefore construct its \textbf{actual} raw metric
measure

\[\nu(u)=(\theta+dd^cu)^n\]

and raw energy E\_theta(u)=(theta,u)\^{}\{n+1\}/(n+1), with zero
reference height. The measure is positive, of mass V, and its finite
model coefficients include ord\_F(X\_0) times the actual intersection
degree on F. Multilinearity identifies it with the epsilon-down-to-zero
limit of actual positive ample measures for L+epsilon H, independently
of H. This identification uses actual metric operators before any limit
and defines no arbitrary inverse measure. These bounded facts are
reconstructed in proposal132 and follow from the complete named
published input units just specified.

Normalize u in PSH(theta) by sup u=0. Its value at the generic trivial
valuation is then zero, as follows already by taking limits of the
defining FS weights. Let u\_j=max(u,-j), and define

\[E_\theta(u)=\inf_{j\geq1}E_\theta(u_j),\qquad
\mathcal E^1_\theta=\{u\in\operatorname{PSH}(\theta):E_\theta(u)>-\infty\}.\]

This is the \textbf{entire actual floor-defined cone}, including
potentials outside E\textasciitilde1. Finite energy is the original
theta-energy condition, with raw normalization. For u in that cone set
C\_u=-E\_theta(u)\textless infinity.

\subsubsection{2. Borel locality at the nef
boundary}\label{borel-locality-at-the-nef-boundary}

We need locality on the strict floor sets, not merely on ordinary open
subsets. For bounded theta-psh a,b, both are finite-energy theta+epsilon
h-psh. The complete actual ample Theorem7.40 yields equality of the two
actual ample metric measures of max(a,b) and a on the Borel set
\{a\textgreater b\}. Multiply by the raw positive ample volume to remove
probability normalization.

Each of these raw measures is a finite polynomial in epsilon as a signed
Radon measure, by the published product construction. Restriction to the
fixed Borel set is a linear operation on these finite signed measures.
Thus their restricted difference is a polynomial identically zero for
every positive rational epsilon. Evaluating against any Borel set, or
solving the finite Vandermonde system at distinct rational epsilon
values, makes every measure coefficient zero. Its constant coefficient
is the actual bounded theta-product. Hence

\[\mathbf1_{\{a>b\}}\nu(\max(a,b))=
\mathbf1_{\{a>b\}}\nu(a).\]

There is no passage of a merely weak limit through restriction to a
discontinuous indicator. This exact coefficient argument retains the
general Borel strict-wins statement of Theorem7.40, whose complete
printed proof was read at literal LF4531--4533 and4565--4586, with
Lemma7.43 LF4552--4564. The source's printed proof uses its
model-component locality and then the actual weighted finite-energy
passage. Its apparent reference to Proposition7.19(iii) at the end is a
source cross-reference issue; the substantive weighted passage is
Proposition7.21(ii), not the scaling assertion in (iii).

\subsubsection{3. Construct the full operator from actual
floors}\label{construct-the-full-operator-from-actual-floors}

Let A\_j=\{u\textgreater-j\}, nu\_j=nu(u\_j), and
sigma\_j=1\_\{A\_j\}nu\_j. For m\textgreater=j, max(u\_m,-j)=u\_j and
\{u\_m\textgreater-j\}=A\_j. Section2 gives
nu\_m\textbar A\_j=nu\_j\textbar A\_j. Since A\_j increases, the actual
sigma\_j form an increasing family of positive Radon measures. Borel
restrictions of finite regular Radon measures on the compact original
space are again regular Radon measures.

The bounded energy identity is the actual mixed formula

\[E_\theta(u_j)=\frac1{n+1}\sum_{r=0}^n
\int u_j\,(\theta+dd^cu_j)^r\wedge\theta^{n-r}.\]

Every mixed measure is positive. Since u\_j\textless=0, this implies

\[\int(-u_j)\nu_j\leq(n+1)(-E_\theta(u_j))\leq(n+1)C_u.\]

On A\_j\^{}c, u\_j=-j, so the mass of rho\_j=1\_\{A\_j\^{}c\}nu\_j is at
most (n+1)C\_u/j. Consequently sigma\_j has mass tending to V. Its
increasing setwise limit is a positive Radon measure sigma of mass V:
the difference sigma-sigma\_j has mass tending to zero, giving
regularity by finite-mass tail approximation. Moreover

\[\|\nu_j-\sigma\|_{\rm TV}\leq
\|\rho_j\|+\|\sigma-\sigma_j\|
\leq\frac{2(n+1)C_u}{j}.\]

Define M\_theta(u)=sigma. This is a measure obtained from the original
actual positive metric floor products and their strict-wins
restrictions. It has the required raw top mass V, gives no mass to
\{u=-infinity\}, and has total-variation convergence of the actual full
floors. For bounded u the construction agrees exactly with the published
bounded operator, since eventually A\_j is all of X\^{}an. All auxiliary
ample choices have already disappeared at each bounded stage.

The same construction with arbitrary real floor levels is increasing and
cofinal with the integers. Changing u by a constant merely shifts the
real levels and leaves curvature products unchanged. Thus the operator
and finite-energy membership are intrinsic under the stated constant
gauge. If V=0, all bounded positive products vanish and all bounded
energies vanish by their constant bounds. The resulting operator is zero
on the entire genuine cone; this does not supply uniqueness modulo
constants or a positive lower-degree replacement.

\subsubsection{4. Every original E1(theta) potential is
integrable}\label{every-original-e1theta-potential-is-integrable}

For bounded normalized theta-psh a,b, monotonicity of the published nef
pairing makes every (theta,a)\textsuperscript{r(theta,b)}\{n+1-r\}
nonpositive. Its all-zero reference height is zero. The identity

\[\int b\,\nu(a)=(\theta,b)(\theta,a)^n-(\theta,0)(\theta,a)^n\]

therefore gives -integral b nu(a) \textless=
-(theta,b)(theta,a)\^{}n.~Expanding (2theta,a+b)\^{}\{n+1\}, all its
terms are nonpositive and the displayed mixed term occurs n+1 times.
Concavity of bounded energy gives

\[\int(-b)\nu(a)\leq-E_{2\theta}(a+b)
=-2^{n+1}E_\theta((a+b)/2)
\leq2^n(-E_\theta(a)-E_\theta(b)).\]

The concavity here follows from the actual ample concavity and
polynomial passage on these bounded inputs; it is not applied to an
unverified nonample ambient domain.

Apply the estimate to a=u\_j and b=v\_l, where u,v are any normalized
original E1(theta) potentials. It bounds integral(-v\_l)nu\_j by
2\^{}n(C\_u+C\_v), uniformly in j,l. First take j to infinity, using
total-variation convergence with the bounded Borel test v\_l. Then let l
increase by ordinary countable monotone convergence. Thus

\[\int(-v)M_\theta(u)\leq2^n(C_u+C_v)<\infty.\]

In particular the constructed measure gives no mass to the pole set of
\textbf{any} original E1(theta) test. This is an actual integrability
statement on the larger intrinsic domain, not an assumed embedding into
E\textasciitilde1.

\subsubsection{5. Strong weighted limits along the actual floor
sequence}\label{strong-weighted-limits-along-the-actual-floor-sequence}

Suppose V\textgreater0; the zero-mass case is immediate. At bounded
theta-psh inputs, define the normalized I and J functionals by the
actual raw energies divided by V and the actual metric measures divided
by V. The fundamental energy identity, the bounds
I/(n+1)\textless=J\textless=nI/(n+1), and Corollary7.31's Holder
estimate transfer from theta+epsilon h by their finite polynomial and
pairing formulas. The constants remain bounded: by the source's explicit
convention at LF473 they depend only on n, while the positive volumes
tend to V\textgreater0. Every input in this passage is bounded and thus
lies in the actual ample finite-energy domain. No inequality on the
larger unbounded cone is assumed.

Let v\_l\textless=v\_k be two floors of a normalized v in E1(theta),
with l\textgreater=k, and set
delta\_kl=(E\_theta(v\_k)-E\_theta(v\_l))/V. The mixed energy difference
formula makes

\[\int(v_k-v_l)\nu(v_l)/V\leq(n+1)\delta_{kl},\qquad
I_\theta(v_k,v_l)\leq n(n+1)\delta_{kl}.\]

All normalized J-values of v\_k,v\_l,u\_j are bounded by
B=max(C\_v,C\_u)/V. Corollary7.31, applied with reference u\_j, gives

\[\int(v_k-v_l)\nu_j/V\leq\delta_{kl}
+C_n\big(n(n+1)\delta_{kl}\big)^{\alpha_n}B^{1-\alpha_n},
\qquad\alpha_n=2^{-n}.\]

Indeed J\_\{u\_j\}(v\_k)-J\_\{u\_j\}(v\_l) equals -delta\_kl plus that
integral. Put delta\_k=(E\_theta(v\_k)-E\_theta(v))/V. It tends to zero;
delta\_kl\textless=delta\_k. Letting l increase and using integrability
yields a bound integral(v\_k-v)nu\_j \textless=V r\_k, uniformly in j,
with r\_k -\textgreater0 given by the right side. The same bound holds
for M\_theta(u), by first passing j to infinity with bounded v\_k-v\_l,
and then l to infinity.

For fixed k, v\_k nu\_j tends to v\_k M\_theta(u) in total variation.
The two tail bounds consequently give

\[v\,\nu_j\longrightarrow v\,M_\theta(u)
\quad\hbox{in total variation}.\]

The same argument gives v\_j nu\_j -\textgreater{} v M\_theta(u) in
total variation. Thus all fixed original E1(theta) weights, as well as
their simultaneous original floors, have the claimed strong
\textbf{floor-sequence} limit. This statement retains the original
operator and weights before limits; it is stronger than merely checking
numerical mass.

\subsubsection{6. Precise remaining
boundary}\label{precise-remaining-boundary}

This construction supplies the actual entire floor-defined nef E1
operator, full raw mass, gauge and auxiliary-ample independence,
integrability of every intrinsic original E1 test, and the stated
weighted countable-floor limits. It does not establish arbitrary
decreasing nets of E1(theta) operators against every original continuous
test, arbitrary model approximations, or the required full variational
equation and uniqueness. The source's Holder estimates on theta-psh test
differences need not separate all continuous functions when theta is
nonample; replacing the original test space by their span would be a new
observable-carrier restriction. Those gaps remain explicit.

No general nonnef big reference theorem or positive lower-degree nonbig
carrier is constructed. A zero top operator on a nonconstant cone cannot
provide uniqueness modulo constants. Original P16.5 remains open, and
the new deduction requires an independent review before any scoped
acceptance count changes.

\subsubsection{7. Source/application
ledger}\label{sourceapplication-ledger-1}

\begin{itemize}
\tightlist
\item
  BJv4: genuine numerical definitions and joint closure; complete
  Theorems7.14/7.18, Lemma7.15 and Proposition7.19 at the selected units
  recorded for proposal132. These are applied only to the verified
  bounded inputs.
\item
  Actual ample G0 body and its raw concavity/metric-product
  compatibility: the exact whole accepted file named in proposal132,
  with verified SHA256
  \nolinkurl{6cd26d798b63a388bd9a5e6dcd0583a7cd9c7417648efe390f4d4d00d28646b9}. Its
  nonample inverse extension is not a grant.
\item
  BJv4 literal LF4074--4125: complete Proposition7.21 and Lemma7.22
  statements/proofs, used to identify the substantive weighted passage
  in the printed locality proof. They are not applied directly to an
  arbitrary intrinsic E1(theta) potential.
\item
  BJv4 LF4175--4208: complete Proposition7.27 statement/proof.
  LF4247--4277 and4280--4294: complete Lemma7.30 and Corollary7.31
  statements/proofs. LF473: constants convention depends only on
  dimension. These are applied at positive ample epsilon on bounded
  inputs, before the proved nef coefficient limit.
\item
  BJv4 LF4531--4533,4552--4586: complete Theorem7.40 and Lemma7.43
  statements/proofs. The Borel polynomial-coefficient argument in
  Section2 is the own extension to bounded nef inputs.
\item
  Primary URL https://arxiv.org/pdf/1801.08229v4 . The actual floor
  construction and its limits are own proofs. Whole-book and
  recursive-foundation verification are not claimed.
\end{itemize}

\subsection{III. Arbitrary mixed weighted nets and both outer
signs}\label{iii.-arbitrary-mixed-weighted-nets-and-both-outer-signs}

This is an unreviewed own continuation. It combines the actual
bounded-domain reconstruction132, the intrinsic floor construction133
and the independently written native121 Section5. It is not an
independent acceptance of those inputs. Original P16.5 remains open.
Neither active native119 nor native122 receives this later material.

\subsubsection{1. Domain and actual starting
constructions}\label{domain-and-actual-starting-constructions}

Retain an algebraically closed trivially valued characteristic-zero
field of arbitrary cardinality, smooth integral projective X of positive
dimension n, actual nef rational L, theta=c1(L), canonical reference0
and raw V=theta\^{}n.~No semiampleness or bigness is assumed. All
potentials belong to the original genuine numerical single-net cone
PSH(theta). The entire intrinsic E1(theta) is defined by the actual
bounded floors and their raw energy, as in133, including endpoints
outside the published auxiliary-ample signed E\textasciitilde1.

The bounded original products are constructed from actual metric
products with their model central multiplicities before any limit. Their
nef positivity and auxiliary independence use verified bounded
membership in E\textasciitilde1 and actual ample perturbations. Borel
winning-set locality uses exact polynomial coefficients, not weak
convergence through indicators. For normalized u in intrinsic E1(theta),
the original floor operator M(u) is constructed by increasing
winning-set restrictions and satisfies

\[M(u)(X^{an})=V,\qquad
\|M(\max(u,-k))-M(u)\|_{TV}\le 2(n+1)C_u/k,
\quad C_u=-E_\theta(u).\]

It integrates every intrinsic E1(theta) test, with the estimate proved
in133. Its fixed-test and simultaneous literal-floor weighted limits are
in total variation. No nonample equation or arbitrary-model limit is a
starting grant.

\subsubsection{2. Energy and top continuity for the original arbitrary
decreasing
net}\label{energy-and-top-continuity-for-the-original-arbitrary-decreasing-net}

Let u\_i decrease pointwise on all of X\^{}an to a genuine generically
finite u. All suprema are bounded on a tail, because they are values at
the generic trivial valuation. At fixed k, max(u\_i,-k) decreases to
max(u,-k) in the verified bounded domain. Bounded energy continuity
gives convergence of the floor energies. Monotonicity yields

\[E_\theta(u)\le\inf_i E_\theta(u_i)
\le E_\theta(\max(u,-k))\quad\hbox{for every }k.\]

Taking k to infinity proves E\_theta(u\_i) decreases to E\_theta(u),
including the extended value -infinity. Simultaneous floors then give
concavity on the entire intrinsic E1 cone; constants shift raw energy by
Vc. The cone is upward finite-energy closed.

If the endpoint is intrinsic E1, choose a common constant B so that
u\_i-B\textless=0 on a tail. Their negative raw energies are bounded by
C=-E\_theta(u-B). The actual floor total-variation error is consequently
at most 2(n+1)C/k uniformly over that original tail and its endpoint.
For each original continuous f, first choose k to make this error small,
and then choose a late original index using bounded-floor weak
continuity. Thus

\[M(u_i)\longrightarrow M(u)\]

weakly against every f in C(X\^{}an). This retains the original
arbitrary-cardinality net. It uses no replacement of the continuous test
space by theta-psh differences and no countable subsequence of the
original net. A sequence is used only for the independent cofinal
reference floors.

This is the useful new observation in native121 Section5.4, now combined
with133: a uniform total-variation floor estimate resolves the earlier
concern about insufficient theta-test separation.

\subsubsection{3. Transfer the energy inequalities to intrinsic
E1}\label{transfer-the-energy-inequalities-to-intrinsic-e1}

Assume V\textgreater0 and divide the actual energies and measures by V
only in the auxiliary normalized I,J notation. Constants do not affect I
or relative J. Every bounded theta-psh input belongs to actual ample
E1(theta+epsilon h); published Proposition7.27 and
Lemma7.30/Corollary7.31 pass to theta by the actual finite-polynomial
coefficient limit. Their dimension-only constants stay bounded as the
positive ample volumes approach V.

For any finite tuple of intrinsic E1(theta) potentials, take their
simultaneous literal floors. The energy terms converge by definition.
Every weighted integral in the I,J identities and Holder bounds
converges by133, applied separately to each original fixed potential and
its actual floors. Hence all those identities and inequalities extend to
this larger intrinsic domain, without placing its endpoints in
E\textasciitilde1. In particular, for a\textgreater=b in E1(theta), with
delta=(E\_theta(a)-E\_theta(b))/V,

\[I(a,b)\le n(n+1)\delta.\]

One way to see the bound before passage is the positive mixed
energy-difference identity: integral(a-b)M(b)/V\textless= (n+1)delta,
whence J\_b(a)\textless=n delta and the fundamental I,J comparison gives
the displayed bound. This reasoning needs order but does not require
equal suprema.

If a\_i decreases to a in E1(theta), Section2 gives
delta\_i-\textgreater0 and therefore I(a\_i,a)-\textgreater0. Normalized
J-values are bounded on a tail because sup a\_i and E\_theta(a\_i)
converge. The transferred Lemma7.30 thus controls integrals of any fixed
intrinsic theta-psh test difference against M(a\_i)-M(a), as well as the
transferred Corollary7.31's relative-J differences.

\subsubsection{4. Fixed and decreasing original E1 weights: full
weighted weak
convergence}\label{fixed-and-decreasing-original-e1-weights-full-weighted-weak-convergence}

Let u\_i decrease to u in E1(theta), mu\_i=M(u\_i), mu=M(u), and let v
be a fixed intrinsic E1(theta) weight. Subtract a fixed constant so
v\textless=0; write v\_k=max(v,-k). The transferred Lemma7.30 applied to
v\_k-0 and u\_i-u gives

\[\int v_k\,d\mu_i\longrightarrow\int v_k\,d\mu\]

for every fixed k. Section2 gives mu\_i-\textgreater mu weakly on the
full original continuous test space. Since -v\_k is a bounded
nonnegative lsc function, the actual primary Lemma7.22, with a constant
net of weights, turns those two facts into weak convergence of
(-v\_k)mu\_i to (-v\_k)mu. This step handles a bounded discontinuous
v\_k; weak convergence alone would not.

The same Corollary7.31 argument as in133 supplies uniform weight tails.
All normalized J-values of u\_i on the tail are bounded by a common B,
as are those of v and its floors. Set

\[\delta_k=(E_\theta(v_k)-E_\theta(v))/V\downarrow0.\]

For l\textgreater=k, apply the bounded or transferred relative-J
inequality to v\_k,v\_l with reference u\_i, and use the ordered I
bound. Passing l to infinity by ordinary countable monotone convergence
gives

\[\int(v_k-v)\,d\mu_i
\le V\left(\delta_k+C_n[n(n+1)\delta_k]^{2^{-n}}B^{1-2^{-n}}\right)=:Vr_k,\]

uniformly on the tail, with r\_k-\textgreater0. The same inequality
holds for mu. For a continuous test f, the error replacing v by v\_k is
at most \textbar\textbar f\textbar\textbar{} Vr\_k at either end. First
choose k and then the original index. Hence v mu\_i converges weakly to
v mu against every original continuous f.

Now also let v\_i decrease to v in intrinsic E1(theta), indexed by the
same arbitrary directed set (a common refinement handles independently
indexed nets). Their J-values are bounded and their raw energy
difference tends to zero. The identical relative-J estimate applied to
v\_i,v, with reference u\_i, gives

\[\int(v_i-v)\,d\mu_i\longrightarrow0.\]

The integrand is nonnegative; thus the total-variation norm of
(v\_i-v)mu\_i tends to zero. Combining this with the fixed-v conclusion
proves

\[v_i M(u_i)\longrightarrow vM(u)\]

weakly as signed Radon measures on the original carrier. Constants are
restored by unweighted weak continuity. Only ordinary countable monotone
convergence was used for fixed-endpoint reference floors. No
arbitrary-Borel-net monotone convergence is substituted for the actual
weighted argument.

\subsubsection{5. Entire mixed products in one intrinsic nef
class}\label{entire-mixed-products-in-one-intrinsic-nef-class}

For u\_1,\ldots,u\_n in E1(theta), their every nonempty subset average
a\_I=\textbar I\textbar\^{}\{-1\}sum\_\{i in I\}u\_i belongs to
E1(theta), by genuine convexity and raw energy concavity. Define the
mixed limit by polarization of the already constructed top operators:

\[M(u_1,\ldots,u_n)=\frac1{n!}
\sum_{\varnothing\ne I\subset\{1,\ldots,n\}}
(-1)^{n-|I|}|I|^nM(a_I).\]

This formula identifies an actual product limit; it does not choose a
measure in response to a target. Indeed for bounded literal floors
u\_\{r,k\}, it is exactly the finite multilinear identity for the actual
original bounded products. Each subset average of those floors decreases
to a\_I with intrinsic E1 endpoint. Section2 therefore shows that the
actual positive mixed floor measures converge to the displayed
expression. The limit is positive, has raw mass V, and is independent of
common cofinal floor levels by the same argument.

We also identify its usual increasing winning-set construction. For
bounded mixed entries, mixed locality follows from bounded top locality
applied to a+sum(lambda\_r c\_r) and b+sum(lambda\_r c\_r) on the fixed
Borel set \{a\textgreater b\}. The class is (1+sum lambda\_r)theta.
Positive rational lambda values and the finite Vandermonde coefficient
argument recover each mixed coefficient. Apply this one entry at a time
on W\_k=\{all u\_r\textgreater-k\}: later actual mixed floors have the
same restriction there.

The positive mixed measure of the n floors is bounded above by n\^{}n/n!
times the top measure of their average, by its positive multinomial
expansion. The normalized floor energies of the entries and their
average are uniformly bounded. Section133's bounded all-E1 integrability
estimate, applied to u\_\{r,k\} and that average, gives a uniform bound
for integral(-u\_\{r,k\}) against the mixed floor measure. Therefore its
mass outside W\_k is O(1/k), by summing the n deep-level events. The
increasing winning-set limit has full mass V and the full mixed floor
measures converge to it in total variation, by the same complement
estimate as Section1. It consequently equals the polarization limit. The
same domination and intrinsic integrability show that all original
E1(theta) tests are integrable against this mixed operator, with no mass
on any such test's poles.

For arbitrary simultaneous decreasing E1(theta) nets
u\_\{r,i\}-\textgreater u\_r, all subset averages decrease with E1
endpoints. Top continuity proves mixed weak continuity by the same
finite polarization. If v\_i decreases to v in E1(theta), Section4
applies to each subset average and gives

\[v_i M(u_{1,i},\ldots,u_{n,i})
\longrightarrow vM(u_1,\ldots,u_n).\]

Thus the integrable weighted statement retains the whole intrinsic floor
domain, all original continuous tests and the original arbitrary nets.
The assertion here concerns one common nef theta; mixed products across
unrelated intrinsic nonample classes are not silently supplied.

If V=0, positive bounded mixed products in this common class vanish, as
do every floor energy, top and mixed operator and weighted measure. The
genuine cone may be nonconstant; this says nothing about uniqueness
modulo constants or a positive lower-degree theory.

\subsubsection{6. Both outer derivatives at every intrinsic E1
potential}\label{both-outer-derivatives-at-every-intrinsic-e1-potential}

For continuous h, let F\_theta(h)=E\_theta(P\_theta h). Proposal132
reconstructs genuine bounded nef envelope membership and the supporting
inequality, by passing actual ample supporting inequalities through
bounded operator/energy limits:

\[F_\theta(h+t f)-F_\theta(h)
\le t\int f\,M(P_\theta h),\]

and its counterpart obtained by interchanging the two obstacles. This is
valid for every continuous original f, both signs of t and every
continuous h. It does not assume that the bounded nef envelope is
continuous.

Let u be any intrinsic E1(theta) potential, and take continuous
majorants h\_i decreasing pointwise to u on compact X\^{}an. Such a
directed net exists for every upper-bounded usc function; finite minima
make the family directed. The envelopes p\_i=P\_theta h\_i decrease to
u: their genuine endpoint lies between u and inf h\_i=u. Energy and top
continuity from Section2 give E\_theta(p\_i)-\textgreater E\_theta(u)
and M(p\_i)-\textgreater M(u).

For a fixed real t, q\_i=P\_theta(h\_i+t f) decreases to
q\_t=P\_theta(u+t f). Genuine nef decreasing closure and all-point
obstacle inequalities prove equality with this envelope. Its lower
competitor u-\textbar t\textbar\textbar\textbar f\textbar\textbar{} and
its upper bound u+\textbar t\textbar\textbar\textbar f\textbar\textbar{}
show genuine intrinsic E1 membership. Apply the bounded supporting
inequalities to h\_i,h\_i+t f and pass the original majorant net using
energy and full-C0 operator continuity. For t\textgreater0 this gives

\[\int fM(q_t)\le
\frac{E_\theta(q_t)-E_\theta(u)}t
\le\int fM(u).\]

For t\textless0 the two integral bounds reverse after division. No
continuity of the envelopes themselves was used.

Finally M(q\_t)-\textgreater M(u) weakly against every original
continuous test as t-\textgreater0. In fact
\textbar q\_t-u\textbar\textless=\textbar t\textbar\textbar\textbar f\textbar\textbar{}
wherever u is finite, and both have the same infinite locus. At every
fixed k, the literal floors obey the global finite bound

\[\|\max(q_t,-k)-\max(u,-k)\|_\infty\le |t|\|f\|_\infty.\]

Proposal132's bounded uniform operator continuity proves convergence of
their bounded floor measures. The suprema and normalized negative raw
energies of q\_t are bounded uniformly for small t, by the two order
bounds. Section1's total-variation floor estimate is therefore uniform
in t. First choose k, then let t-\textgreater0. This proves the asserted
full-C0 weak continuity. The two supporting inequalities now squeeze
both signs to

\[\left.\frac d{dt}\right|_{t=0}E_\theta(P_\theta(u+t f))
=\int f\,M(u).\]

This is also the actual outer energy derivative: for an upper-bounded
usc obstacle admitting a genuine intrinsic E1 minorant, the envelopes of
decreasing continuous majorants decrease to its genuine envelope, and
Section2 gives energy convergence. Thus the infimum defining the outer
energy equals the envelope energy used above. At V=0 both sides vanish.

\subsubsection{7. Source applications and remaining original
obligations}\label{source-applications-and-remaining-original-obligations}

The whole final own133 and the whole native121 proof were read before
this continuation. Exact input grants remain confined to the genuine
actual-ample G0 body. The bounded facts are the unreviewed132
construction under BJv4 Theorems7.14/7.18, Lemma7.15, Proposition7.19
and complete nef envelope Theorem5.20; Lemma5.17 is only used on actual
ample classes. Borel locality uses complete actual ample Theorem7.40 and
Lemma7.43, transferred by an explicit finite coefficient argument.

The complete primary Lemma7.22 at literal LF4082--4104 is used only
after unweighted weak convergence and convergence of the integrals of a
nonnegative lsc weight have separately been proved. Proposition7.27 at
LF4175--4208 and complete Lemma7.30/Corollary7.31 at LF4247--4294 are
applied first to bounded actual ample inputs, then to the verified nef
coefficients and finally through proved intrinsic weighted floor limits.
The dimension-only constant convention was read at LF473. Primary URL:
https://arxiv.org/pdf/1801.08229v4 . Whole-book and recursive-foundation
verification are not claimed.

The current result is an unreviewed entire intrinsic nef
operator/mixed/weighted/differentiability theorem. It does not prove
divisorial weak energy upper semicontinuity and target-integral
continuity on energy sublevels, coercivity and actual variational
attainment, domination or uniqueness for every nef big class. It
constructs no general big nonnef minimal/model reference and no general
nonbig positive lower-degree carrier. Original P16.5 remains open; no
solved-original or independently accepted scoped count changes.

\subsection{IV. Weak energy, intrinsic targets and
coercivity}\label{iv.-weak-energy-intrinsic-targets-and-coercivity}

This is a new unreviewed own continuation of132--134, not an independent
acceptance or complete original P16.5 solution. Active native122 is
frozen and receives none of this material. The actual operator and
intrinsic energy domain below are the original floor constructions; no
auxiliary ample E1 membership of unbounded endpoints is assumed.

\subsubsection{1. Retain the whole intrinsic
domain}\label{retain-the-whole-intrinsic-domain}

Let k,X,L,theta,n and V=theta\^{}n be as in134: arbitrary-cardinality
algebraically closed trivially valued characteristic-zero field, smooth
integral projective X of positive dimension, and any actual nef rational
L, canonical reference0. For the positive-mass assertions assume
V\textgreater0. Every genuine potential is normalized by sup
u=u(v\_triv)=0. Write C(u)=-E\_theta(u)\textgreater=0. All original
E1(theta) functions, including those outside E\textasciitilde1, are
retained.

Use the verified-domain bounded metric pairing and its nef monotonicity
from the primary Theorems7.14/7.18 and Lemma7.15. Genuine normalized nef
compactness and the envelope property come from complete primary
Theorem5.20 and the genuine-cone constructions132/native121. Actual
ample G0 is used solely for actual ample auxiliary classes and their
energy upper semicontinuity.

\subsubsection{2. Divisorial weak upper semicontinuity of the intrinsic
energy}\label{divisorial-weak-upper-semicontinuity-of-the-intrinsic-energy}

Fix an actual ample H, h=c1(H), and an integer k. On genuine theta-psh
scalars with -k\textless=b\textless=0, the published bounded height
pairing gives

\[E_{\theta+\epsilon h}(b)
=\frac1{n+1}\sum_{l=0}^{n+1}\binom{n+1}l\epsilon^l
(\theta,b)^{n+1-l}(h,0)^l.\]

The all-zero reference height is zero. For 1\textless=l\textless=n,
monotonicity replacing each of the m=n+1-l theta-slots between the
constant metrics -k and0 gives

\[-mk\,\theta^{n-l}h^l
\le(\theta,b)^m(h,0)^l\le0.\]

The l=n+1 term is zero. Consequently

\[|E_{\theta+\epsilon h}(b)-E_\theta(b)|
\le\frac{k}{n+1}\sum_{l=1}^n
\binom{n+1}l(n+1-l)\epsilon^l\theta^{n-l}h^l.\]

This tends to zero uniformly over the original bounded domain as
rational epsilon decreases to zero. Each energy on the left auxiliary
actual ample class is upper semicontinuous in divisorial topology by G0;
the inclusion of the genuine theta-cone into that ample cone is
continuous in the same coordinates. A uniform limit of upper
semicontinuous real functions is upper semicontinuous. Thus bounded
E\_theta is weakly upper semicontinuous there.

The genuine floor map u-\textgreater max(u,-k) is continuous in
divisorial topology because it is coordinatewise maximum with a
constant, and genuine endpoints are all-point determined. On the
normalized cone its values lie between -k and0. Hence

\[E_\theta(u)=\inf_{k\ge1}E_\theta(\max(u,-k))\]

is an infimum of upper semicontinuous functions, so is upper
semicontinuous with extended values on the entire normalized genuine
cone. Therefore \{u:sup u=0,C(u)\textless=B\} is a closed subset of the
genuinely compact normalized nef cone, and is weakly compact for every
finite B. This is compactness of actual original functions, not a
completion or a countability assertion.

\subsubsection{3. A quadratic contraction estimate from the actual
metric
products}\label{a-quadratic-contraction-estimate-from-the-actual-metric-products}

For normalized bounded u, set

\[A_l=\int(-u)(\theta+dd^cu)^l\wedge\theta^{n-l},\quad0\le l\le n.\]

The measures are actual positive bounded mixed products. Since
theta\^{}n at the canonical zero metric is V delta\_vtriv and
u(v\_triv)=0, A\_0=0. The actual energy identity says sum
A\_l=(n+1)C(u).

For 0\textless a\textless=1, the genuine convex combination au has
curvature (1-a)theta+a(theta+dd\^{}cu). Its actual mixed expansion gives

\[C(au)=\frac1{n+1}\sum_{r=0}^n\sum_{l=1}^r
\binom rl a^{l+1}(1-a)^{r-l}A_l.\]

Every term is nonnegative. Using a\textsuperscript{\{l+1\}\textless=a}2
and sum\_\{r=l\}\^{}n binom(r,l)=binom(n+1,l+1)\textless=2\^{}\{n+1\},
obtain

\[C(au)\le K_n a^2 C(u),\qquad K_n=2^{n+1}.\]

For every original intrinsic E1 u, apply this to its actual normalized
floors. Their scaled metrics decrease to au and their thresholds aj are
cofinal real reference thresholds. Intrinsic decreasing energy
continuity from134 therefore gives the same bound on the entire floor
domain. This is not an auxiliary-ample estimate on an unverified
unbounded domain.

\subsubsection{4. Exact equivalence of the intrinsic dual finiteness
condition and
integrability}\label{exact-equivalence-of-the-intrinsic-dual-finiteness-condition-and-integrability}

Let mu be a positive original-carrier Radon measure of raw mass V.
Define

\[D_\theta(\mu)=\sup_{u\in E^1(\theta)}
\left(E_\theta(u)-\int u\,d\mu\right).\]

Raw mass V makes this objective invariant under constants. If D is
finite, every genuine intrinsic E1 potential is integrable: after
normalizing u\textless=0 it is bounded above, and integral(-u)=infinity
would make the objective infinite.

Conversely suppose every intrinsic E1(theta) potential is mu-integrable.
The integrals of -u are uniformly bounded on the normalized set
C(u)\textless=1. Otherwise choose u\_j in that set with
integral(-u\_j)\textgreater=2\^{}j. Let s\_m=sum\_\{j=1\}\^{}m
2\^{}\{-j\}u\_j, retaining the remaining convex weight on the genuine0
potential. These actual convex combinations decrease, have generic
value0 and raw energy at least -1 by concavity. Genuine decreasing
closure gives their actual endpoint s, and energy continuity makes s
intrinsic E1. Ordinary countable monotone convergence gives
integral(-s)\textgreater=sum\_j1=infinity, a contradiction. Denote the
finite uniform bound by B\_mu.

For arbitrary normalized u with C=C(u)\textgreater1, choose a=(K\_n
C)\^{}\{-1/2\}. Section3 gives C(au)\textless=1, whence

\[\int(-u)\,d\mu\le B_\mu\sqrt{K_n C}.\]

For C\textless=1 the same integral is at most B\_mu. Thus the objective
-C+integral(-u) is uniformly bounded above, proving
D\_theta(mu)\textless infinity. Its exact intrinsic finiteness condition
is therefore equivalent to integrability of every original intrinsic E1
test. This statement characterizes the condition, not yet the complete
image of the operator.

Every actual top and common-class mixed output constructed in133/134
integrates every such test and has raw mass V. Hence its dual energy is
finite. This direction follows from the already constructed metric
operator, not from an assumed inverse equation.

\subsubsection{5. Coercivity and the remaining attainment
issue}\label{coercivity-and-the-remaining-attainment-issue}

If D=D\_theta(mu)\textless infinity, choose once for all
0\textless a\textless=1/(2K\_n). Applying the dual inequality to the
genuine au and the quadratic estimate gives

\[a\int(-u)\,d\mu\le D+C(au)
\le D+K_n a^2 C(u).\]

Therefore the normalized objective satisfies

\[E_\theta(u)-\int u\,d\mu
\le-\frac12 C(u)+D/a.\]

Every maximizing sequence or net with objective tending to finite D has
uniformly bounded original C-values on a tail. Section2 gives an actual
weakly convergent subnet with intrinsic E1 endpoint. The same inequality
and the integrability argument preserve the exact raw original target
space.

However, upper semicontinuity of the energy alone does not prove upper
semicontinuity of this objective: target integration must be shown
continuous, or a separate strong-Cauchy construction must be completed.
The present argument does not supply that missing passage, actual
maximizer attainment, or uniqueness. General nef big variational
equations therefore remain open here.

\subsubsection{6. Precise source use and next mathematical
gate}\label{precise-source-use-and-next-mathematical-gate}

The bounded height formula, monotonicity and actual mixed operators use
the complete BJv4 Theorems7.14/7.18, Lemma7.15 and Proposition7.19 at
the verified bounded domain. The genuine nef envelope hypothesis is
complete Theorem5.20, not a relabelled ample theorem. The exact G0 body
grants auxiliary actual ample energy usc; the original nef weak usc is
the own uniform coefficient argument above. The whole final own133/134
and native121 Section5 supply only the actual intrinsic floor products,
positivity, concavity and decreasing energy passage reconstructed in
those drafts.

The convex-series idea is compared with the complete printed BJv4
Theorem9.20 beginning at literal LF5266, but the argument above is
independently given on the larger intrinsic nef domain, with a different
explicit quadratic scaling proof. No general nonample result is quoted
from that ample theorem. Primary URL: https://arxiv.org/pdf/1801.08229v4
. No whole-book or recursive foundational proof verification is claimed.

The useful next problem is target-integral continuity on the original
compact energy sublevels, or a complete intrinsic I-Cauchy maximizer
construction, followed by all-point uniqueness for nef big L. Arbitrary
theta-psh test differences need not span all original continuous
functions. Any use of the ample source's PL-theta generation or
divisorial finite-energy Dirac tests must reconstruct their nonample
analogues. Replacing the carrier or assuming unbounded auxiliary-ample
E1 membership would not resolve that gate. General big nonnef references
and general nonbig nonzero numerical data are still not constructed.
Original P16.5 remains open and every assertion in this note requires
independent review.

\subsection{V. Complete independent reconstruction and final
scope}\label{v.-complete-independent-reconstruction-and-final-scope}

PASS for proposals132--135 in the stated nef domain, relative solely to
the inline actual-ample G0 grant. This covers arbitrary-cardinality k,
smooth integral projective X, n≥1, every actual nef rational L and
canonical reference0, and the entire intrinsic floor E1(θ). It does not
accept unchanged original P16.5, an inverse theorem, attainment or
uniqueness. The urn passage documents the ample antecedent used in
independently reconstructing the nef conclusion associated with
premise6; it grants no nef assertion.

Genuine membership was reconstructed as PSH(θ)=∩ε\textgreater0
PSH(θ+εh). Forward section embeddings preserve the original scalar
because generating canonical section norms have maximum0. For the
reverse inclusion, shifted finite continuous ample FS approximants form
one directed decreasing net: strict comparison with a finite continuous
old minimum gives an open cover, and compactness gives a common later
term. This handles poles and avoids the defective general
extended-valued flattening lemma. Closed intersections inside a
normalized ample cone give genuine compactness and the envelope
property. Lemma5.17 remains ample-only. No unbounded intrinsic endpoint
is placed in E\textasciitilde1; Remark7.16's distinction is retained.

Bounded inputs belong to a common actual ample E1 space. Actual tensor
identities L=(L+H)−H construct signed metric products before limits,
with canonical references and continuous dpsh compatibility. At model
stages the coefficient is b\_F∏r(θr,𝓧+D\_r)\textbar F, with
b\_F=ord\_F(t) and v\_F=b\_F\^{}−1 ord\_F\textbar k(X). The collar
identity retains its fixed reference term. Multilinearity proves
auxiliary independence; positive actual ample perturbations give nef
positivity and raw mass V=θ\^{}n.~Borel strict-winning locality follows
by restricting a finite measure-valued polynomial and extracting its
exact coefficients. It does not pass weak convergence through
indicators.

For normalized \(u\), put \(C_u=-E_\theta(u)\). Positivity of the bounded mixed energy
formula gives \(\int(-u_k)M(u_k)\leq(n+1)C_u\). The increasing strict-winning
restrictions therefore have full mass \(V\) and yield
\(\|M(u_k)-M(u)\|_{\mathrm{TV}}\leq2(n+1)C_u/k\). Their
limit is Radon and agrees with the bounded operator. Real cofinal floors
and constant shifts preserve it. Expanding the bounded nef height gives
∫(−v)M(u)≤2\^{}n(C\_u+C\_v), first by bounded-test TV convergence, then
countable monotone convergence. Thus every intrinsic E1 test is
integrable, including endpoints outside E\textasciitilde1.

The normalized I/J and Hölder inequalities are applied first at positive
ample ε to bounded actual inputs. Polynomial pairing limits preserve
them with dimension-only constants as Vε→V\textgreater0. Ordered floors
satisfy I(a,b)≤n(n+1)(Eθ(a)−Eθ(b))/V. The relative-J estimate gives
uniform vanishing weight tails; bounded-weight TV convergence then
proves fixed and simultaneous intrinsic weighted floor limits. Only
after these limits are established are the inequalities transferred to
intrinsic E1.

For an original decreasing net with E1 endpoint, floor energy continuity
proves full decreasing energy continuity. A common upper shift bounds
all tail costs, making the TV floor error uniform over the original
directed set. Choosing the floor first and the original index second
proves weak convergence on the full original C(X\^{}an). For bounded
discontinuous floors, integral convergence is separately obtained from
the transferred inequality before applying Lemma7.22. Uniform weight
tails handle intrinsic weights and simultaneous decreasing weights. More
generally, for a genuine decreasing psh weight integrable at the
endpoint, polarized intrinsic heights give the lower integral bound;
fixed floors give the upper bound. The resulting integral convergence
establishes Lemma7.22's antecedent before weighted weak convergence. No
arbitrary-Borel-net monotone convergence or restricted test carrier is
used.

Finite subset-average polarization identifies common-θ mixed products as
limits of actual positive bounded mixed floors. Mixed coefficient
locality fixes their restrictions on W\_k=\{all entries\textgreater−k\}.
Positive multinomial domination by the average's top product bounds the
complementary mass by O(1/k). Thus their increasing winning-set
construction and polarization limit coincide, with full mass, TV floor
convergence, every intrinsic test integrable, and arbitrary decreasing
mixed weighted-net convergence.

Ample supporting inequalities pass through the bounded nef limits.
Decreasing continuous majorants of any intrinsic E1 base then give
supporting inequalities at that base. The perturbed envelope lies
between u±\textbar t\textbar\textbar\textbar f\textbar\textbar; its
fixed floors converge uniformly and its floor TV errors are uniform for
small t. Both quotient signs are consequently squeezed to ∫f M(u). The
same majorant argument identifies envelope energy with outer energy.

The bounded ample-height coefficient estimate is uniform on −k≤b≤0.
Hence bounded nef energy is weakly usc; continuous divisorial floor maps
and an infimum establish intrinsic usc. The stated normalized cost sets
\{−Eθ≤B\} are closed in the genuine compact cone. Positive mixed
expansion, with A0=0 at the canonical generic trivial valuation, proves
C(au)≤2\^{}(n+1)a²C(u); cofinal scaled floors extend it intrinsically.
Actual decreasing convex series prove that integrability of all
intrinsic tests bounds their integrals uniformly on C≤1. Scaling then
proves dual finiteness, and every top/common-class mixed output has
finite dual energy. For raw mass V\textgreater0, the objective obeys
Eθ(u)−∫u dμ≤−C(u)/2+2\^{}(n+2)Dθ(μ). This proves coercivity without
asserting objective usc or attainment. At V=0 all positive common-class
products, floor energies and weighted measures vanish; no probability
normalization or uniqueness is inferred.

Fresh complete pertinent primary units were checked: BJ v4, March2022,
printed pp21,25--26,36--38,42--47,57--63,79--87,91--92; BE v3,
January2021, pp50,55--58,65--66; Fujino--Miyamoto v1 Theorem1.4 and its
pp4--5 proof; Stacks01Q3's complete proposition/proof. Relevant
published BJ units were matched independently. Harmless corrections are
the reversed pairing sign in 7.21(ii), the 7.40 limiting reference to
7.21(ii), and algebraic indices/degrees in 3.31/3.33/7.21/7.30; the n=1
Hölder case uses direct Cauchy--Schwarz. The abstract flattening failure
is substantial, but the supplied finite-continuous nef argument
independently replaces that application. No historical verdict, hash,
nonample note or other campaign supplies a premise.

\section{Appendix C. Genuine polar tests and original continuous
observables}\label{appendix-c.-genuine-polar-tests-and-original-continuous-observables}

The following complete mathematical derivation is retained as supporting
proof, with its earlier scope statements. It is not a historical-verdict
grant. Equation numbers in this appendix are local to it. The main
candidate and the fresh review criteria govern the new combined scope.

\section{Nef polar capacity and control of the original continuous
observables}\label{nef-polar-capacity-and-control-of-the-original-continuous-observables}

This is the unreviewed repair2 of own continuation136. The earlier
conditional PL paragraph wrongly assumed both small signs of a PL
function are ample-psh; Section5 now uses a difference of two ample-psh
functions, without that assumption. The earlier published unreviewed
source and hash remain preserved. This is an unreviewed own continuation
for original P16.5. The entire intrinsic nef constructions in
proposals133--135 are mathematical inputs to be independently
reconstructed, not granted original solutions. No active frozen
native123 or native124 input is changed or supplied hot feedback. This
note proves additional implications of those constructions; it does not
claim variational attainment or uniqueness.

\subsection{1. Domain and normalization}\label{domain-and-normalization}

Let X be smooth integral projective of dimension n\textgreater=1 over an
algebraically closed trivially valued field of characteristic zero and
arbitrary cardinality. Retain an actual nef rational L, its genuine
numerical class theta, canonical zero metric, and positive raw top mass
V=theta\^{}n.~The measures in this note are probability-normalized only
after constructing the actual metric operator:
mu\_u=V\^{}\{-1\}M\_theta(u). Write e(u)=V\^{}\{-1\}E\_theta(u) and
C(u)=-e(u) for a normalized genuine potential u with u(v\_triv)=0, hence
u\textless=0. The canonical zero measure is delta\_v\_triv by the actual
trivial-model component formula, not a new reference measure.

The bounded products and all mixed integration-by-parts identities are
the actual published metric products verified in132. For entire
intrinsic E1(theta),133 supplies the literal floor construction, full
mass, mixed integrability and strict winning locality;134 supplies
weighted limits and the numerical I/J inequalities on that larger
domain. None of these is replaced by an ambient unbounded
E\textasciitilde1 membership assertion.

Define, for Borel B on the original X\^{}an,

\[\operatorname{Cap}_\theta(B)=\sup\{\mu_b(B):b\in\operatorname{PSH}(\theta),\ -1\leq b\leq0\}.\]

The full original carrier and Borel sets are retained. This need not be
a newly asserted Choquet capacity; only the displayed set function and
estimates are needed below.

\subsection{2. A capacity bound for every genuine
pole}\label{a-capacity-bound-for-every-genuine-pole}

Let a be ANY normalized genuine theta-psh function, with no
finite-energy hypothesis. Its bounded floors a\_j=max(a,-j) are
theta-psh. For any bounded theta-psh b between -1 and0, telescoping from
the zero metric gives

\[M_\theta(b)-M_\theta(0)=\sum_{r=0}^{n-1}dd^c b\wedge\theta_b^r\wedge\theta_0^{n-1-r}.\]

Integration by parts against -a\_j turns the r-th integral into the
difference of the integrals of -b against the two positive mixed
products having a\_j-curvature or zero-curvature in the remaining slot.
Both have raw mass V. Since 0\textless=-b\textless=1, the absolute value
of each difference is at most V. The initial integral against
M\_theta(0)=V delta\_v\_triv is zero because a\_j(v\_triv)=0. Hence

\[\int(-a_j)\,d\mu_b\leq n.\]

Monotone convergence along the COUNTABLE literal floors, without
invoking a Borel monotone theorem for arbitrary nets, proves the same
bound for -a. Markov's inequality uniformly over the defining b gives

\[\operatorname{Cap}_\theta\{a\leq-t\}\leq n/t,\qquad t>0.\tag{1}\]

No ample lower bound for theta was used: actual bounded positive mixed
products and equal numerical mixed masses suffice. In particular every
bounded genuine metric measure gives every genuine theta-psh pole set
zero mass.

\subsection{3. An intrinsic finite-energy function with the same
poles}\label{an-intrinsic-finite-energy-function-with-the-same-poles}

For 0\textless alpha\textless1/(n+1), set chi(x)=1-(1-x)\^{}alpha on
x\textless=0. It is increasing convex, vanishes at zero, and has slopes
between0 and1. A precise genuine approximation can be made with
piecewise affine chord majorants on {[}-R\_j,0{]}, constant tails to the
left, and nested subdivisions. Let R\_j increase to infinity, subdivide
every old interval further, and subdivide new intervals finely enough
that the maximal chord error on each fixed compact interval tends to
zero. Convexity makes refinement lower the chord majorant. On a new left
interval its chord is at most the old constant tail; on the old
intervals refinement is at most the old chord. The resulting functions
chi\_j decrease to chi on the extended interval {[}-infinity,0{]}, with
chi(-infinity)=-infinity. Each finite convex piecewise affine chi\_j is
the maximum of finitely many affine functions sx+d with
0\textless=s\textless=1, including a constant branch. Consequently
chi\_j(a) is bounded genuine theta-psh by scaling with the genuine zero
metric, constant translation and finite maxima. The genuine nef cone is
decreasingly closed at the generically finite limit, so c=chi(a) is
genuine theta-psh, normalized at v\_triv.

Here is a full energy criterion for this intrinsic cone. For ANY
normalized genuine c, put c\_s=max(c,-s) and mu\_s=mu\_\{c\_s\}. Actual
bounded energy monotonicity of the mixed terms gives

\[C(c_s)=\frac1{n+1}\sum_{r=0}^{n}\int(-c_s)\,V^{-1}\theta_{c_s}^r\theta_0^{n-r}
\leq\int(-c_s)\,d\mu_s.\]

For completeness the ordering of consecutive mixed integrals follows by
integration by parts: their difference is the nonnegative Dirichlet
pairing of c\_s with itself against the remaining positive bounded
theta-products. It is the bounded numerical I-positivity identity, first
obtained in the actual ample perturbations and then coefficientwise at
epsilon=0. The initial term is zero.

If 1\textless=t\textless=s, strict winning locality identifies mu\_s and
mu\_t on \{c\textgreater-t\}. Their total probability masses are equal,
so their masses on the complementary Borel set \{c\textless=-t\} are
equal. The positive mixed expansion of mu\_\{c\_t/t\} contains
t\^{}\{-n\}mu\_t, since zero is genuinely theta-psh and t\textgreater=1.
This proves

\[\mu_s\{c\leq-t\}=\mu_t\{c\leq-t\}
\leq t^n\operatorname{Cap}_\theta\{c\leq-t\}.\]

Layer cake, using sets with \textgreater= in the threshold which gives
the same Lebesgue integral, now yields

\[C(c_s)\leq1+\int_1^\infty t^n\operatorname{Cap}_\theta\{c\leq-t\}\,dt.\tag{2}\]

No weighted capacity exponent has been discarded. For c=chi(a), solving
chi(a)\textless=-t gives

\[a\leq1-(1+t)^{1/\alpha},\]

so (1) bounds the capacity in (2) by O(t\^{}\{-1/alpha\}). The integral
converges because alpha\textless1/(n+1). Thus the intrinsic floor energy
of c is finite. Its pole set is exactly that of a.

By133, every intrinsic output mu\_u, u in E1(theta), integrates
-c.~By135's intrinsic dual characterization, every target of finite dual
energy integrates -c as well. Therefore both outputs and all intrinsic
finite-dual-energy targets give EVERY genuine theta-psh pole set zero
mass. This conclusion does not claim that every theta-psh a itself has
finite energy.

\subsection{4. Uniform polar control on energy-bounded
families}\label{uniform-polar-control-on-energy-bounded-families}

The intrinsic estimate in133, after division by V, is

\[\int(-v)\,d\mu_u\leq2^n\bigl(C(u)+C(v)\bigr)\]

for normalized intrinsic finite-energy u,v. Apply it to v=c=chi(a). If
C(u)\textless=B, then, for R\textgreater0,

\[\mu_u\{a\leq-R\}\leq
\frac{2^n(B+C(c))}{(1+R)^\alpha-1}.\tag{3}\]

This uniform estimate is on the actual entire space and includes all
nondivisorial points. It gives a useful substitute for an unproved ample
representation of every unbounded intrinsic potential.

\subsection{5. Original PL tests under a retained big-class
decomposition}\label{original-pl-tests-under-a-retained-big-class-decomposition}

Suppose additionally that actual rational bundles and an actual
effective rational divisor have been supplied with L=A+D, where A is
ample. The existence of this decomposition is not granted by this
conditional statement. Let w be its actual canonical divisor metric,
normalized at the generic trivial valuation. On the complement of D, its
formal local curvature is that of the actual ample A metric. Global
genuine membership is verified directly by sections: if m clears the
rational data, multiply a generating family of mA by the actual
canonical section of mD; the resulting possibly singular FS maximum is
w. This checks genuine membership without redefining the cone.

For any original PL test f, the actual ample PL generation theorem
expresses a positive rational multiple t f as b-b\_prime, where b and
b\_prime are bounded continuous genuine A-psh PL functions. Explicitly,
split the finite rational linear combination of A-FS functions into
positive and negative sums, and divide by a rational number at least the
total weight on either side; add the genuine zero A metric for the
unused weight. Constants can be included in one summand. Actual section
multiplication makes w+b and w+b\_prime genuine theta-psh. For
k\textgreater=1 define bounded genuine tests

\[w_k=\max(w+b',-k),\qquad z_k=\max(w+b,-k),\qquad d_k=z_k-w_k.\]

Then \(|d_k|\leq t\|f\|_\infty\)
on the entire original space, including common poles. Set
\(B_f=\max(\|b\|_\infty,\|b'\|_\infty)\).
On \(\{w>-k+B_f\}\), both winning floors are uncut, so \(d_k=tf\).
Consequently

\[|d_k-tf|\leq2t\|f\|_\infty\,
\mathbf1_{\{w\leq-k+B_f\}}.\tag{4}\]

Let u\_i be a normalized intrinsic energy-bounded net which is Cauchy
for the probability-normalized numerical I. For each FIXED k, the
transferred intrinsic Holder estimate of134, with the fixed bounded
genuine tests z\_k and w\_k, gives

\[\int d_k\,(d\mu_{u_i}-d\mu_{u_j})\longrightarrow0.\]

Its other numerical J factors stay bounded on the fixed energy sublevel.
One can choose the zero reference in the polarized estimate:
J(u\_i)=C(u\_i) and J(z\_k),J(w\_k) are fixed finite numbers. This
application does not assume d\_k is continuous, or that w belongs to
E1(theta).

Now use (3) for a=w and (4) to make the replacement of d\_k by tf
uniformly small over the entire net as k increases. First choose k, then
choose sufficiently late i,j. It follows that the integrals of EVERY
original PL f against mu\_\{u\_i\} are Cauchy. Uniform density of PL on
compact X\^{}an and the common unit masses imply the same for EVERY f in
C0(X\^{}an). The resulting bounded positive linear functional has mass
one; by the compact Radon representation theorem it is a Radon
probability measure on the original carrier. Thus mu\_\{u\_i\} has a
unique original weak limit. This is a consequence about the measure
image of an I-Cauchy energy-bounded family, not a proof that the family
itself has an intrinsic E1 limit.

\subsection{6. Exact remaining gate}\label{exact-remaining-gate}

The preceding capacity and polar results hold for every positive-mass
nef L, without semiampleness. The PL-observable conclusion additionally
retains an actual supplied ample-plus-divisor decomposition. They remove
two specific obstacles: poles cannot carry finite-dual-energy mass, and
original continuous observables can be controlled using bounded genuine
tests even when an ample PL direction is unavailable globally in the nef
cone.

They do not show completeness or separation of numerical I on the entire
intrinsic E1(theta), continuity of the target integral on a general
weakly convergent energy sublevel, or attainment of the variational
supremum. A weak limit of the output measures is not identified with the
output of a potential without an additional argument. General nonnef
minimal references and nonbig nonzero carriers remain outside this note.

\subsection{7. Sources and proof
provenance}\label{sources-and-proof-provenance}

The sole equation grant is the exact accepted actual-ample G0 body. Its
approximation, numerical identities and PL density are applied only in
that domain. Bounded nef products are the published BJ v4 pairing and
measure constructions after verifying their bounded inputs in one common
ample E1 cone. BJ v4 Sections7.27,7.30--7.31 and7.40 supply the complete
printed numerical and strict-locality proofs in their original ample
domains;133--134 explicitly transfer their bounded polynomial identities
and intrinsic floor limits rather than assume unbounded ambient
membership. The new capacity deduction in Sections2--4 is reconstructed
here. It also parallels the already accepted ample polar-test argument
in the scoped119 paper Section4.2, with the bounded nef positivity and
raw-mass identities checked explicitly.

Primary source: https://arxiv.org/pdf/1801.08229v4 . No complete
original P16.5 closure or new independent PASS is claimed for this
unreviewed continuation.

\section{Appendix D. Divisorial atomic attainment and all-point
numerical
separation}\label{appendix-d.-divisorial-atomic-attainment-and-all-point-numerical-separation}

The following complete mathematical derivation is retained as supporting
proof, with its earlier scope statements. It is not a historical-verdict
grant. Equation numbers in this appendix are local to it. The main
candidate and the fresh review criteria govern the new combined scope.

\section{Divisorial atomic equations and separation of intrinsic nef
I}\label{divisorial-atomic-equations-and-separation-of-intrinsic-nef-i}

This is a new unreviewed own deduction. It depends on independently
reconstructing the entire intrinsic operator, weighted inequalities,
both outer derivative signs, compact cone, energy usc and coercivity
in133--135. It changes neither original closure counts nor either active
frozen native123/124 packet. It supplies finite divisorial atomic
attainment and all-point I separation, not attainment for every
finite-energy target.

\subsection{1. Literal scope and
measures}\label{literal-scope-and-measures}

Retain the field, smooth integral projective original X,
n\textgreater=1, actual nef rational L, genuine numerical cone,
canonical reference and positive raw V=theta\^{}n of136. Use
e=E\_theta/V and the derived probability operator mu\_u=M\_theta(u)/V.
All equations below are literal raw metric equations after multiplying
by V. The entire intrinsic E1(theta) is defined by actual bounded floor
energies. No ambient unbounded E\textasciitilde1 identification occurs.

Fix any finite probability combination

\[\nu=\sum_{j=1}^m b_j\delta_{v_j},\qquad b_j\geq0,\quad\sum b_j=1,\]

where v\_j are original divisorial valuations, optionally including the
generic trivial valuation. Normalization means u(v\_triv)=0. The weak
topology is the existing pointwise topology on divisorial coordinates of
the genuine cone.

\subsection{2. Actual finite bounds and weak
continuity}\label{actual-finite-bounds-and-weak-continuity}

Choose one actual ample Omega=L+H. Every normalized genuine theta-psh u
is Omega-psh by adding the zero nef H metric. At each divisorial v, the
genuine ample theory gives

\[-T_\Omega(v)\leq u(v)\leq0,\qquad T_\Omega(v)<\infty.\]

For an independent primary check, BJ v4 Proposition1.28, complete
printed LF922--931, proves finite linear growth: realize v=t ord\_F on a
normal projective birational model, intersect the effective divisor
div(pi*s)-ord\_F(s)F with an ample (n-1)-fold power, and obtain
ord\_F(s)\textless=C m for sections of mOmega. Lemma4.46, complete
LF2703--2712, identifies this section bound with the normalized psh
bound by finite maxima and the actual ample approximation theorem. The
generic trivial bound is zero. Thus every u is finite at every v\_j, and

\[0\leq-\int u\,d\nu\leq\sum_j b_j T_\Omega(v_j).\]

This also directly proves D\_theta(nu)\textless infinity, since
e(u)\textless=0. The integral is a finite linear combination of the
coordinates defining the weak topology and hence is continuous on the
ENTIRE normalized genuine cone. No continuity theorem for general
targets is being presumed.

\subsection{3. Actual variational
attainment}\label{actual-variational-attainment}

Let

\[D_\theta(\nu)=\sup_{u\in\mathcal E^1(\theta),\ u(v_{\rm triv})=0}
\left(e(u)-\int u\,d\nu\right).\]

The supremum is finite by Section2, and at least zero by u=0. Choose a
maximizing net. Its eventual energy costs C(u\_i) are bounded: this
follows from135's intrinsic coercivity, or directly here from the
uniform finite bound for -integral u\_i dnu and the eventual lower bound
on the objective. The genuine normalized cone is weakly compact by the
explicit nef intersection proof in121/132, and135 proves intrinsic
energy usc and compactness of each energy sublevel. Pass to a convergent
subnet with genuine limit u and C(u)\textless=B. Therefore u is in the
entire intrinsic E1(theta). Energy usc and the finite-coordinate
continuity give

\[e(u)-\int u\,d\nu\geq\limsup_i\left(e(u_i)-\int u_i\,d\nu\right)=D_\theta(\nu).\]

So the supremum is attained. This argument works with
arbitrary-cardinality nets, without a countable extraction of potential
approximants.

\subsection{4. Verification of the original
equation}\label{verification-of-the-original-equation}

For f in the FULL C0(X\^{}an) and real t, put q\_t=P\_theta(u+t f).
Bracketing by u plus or minus
\textbar t\textbar\textbar\textbar f\textbar\textbar{} shows q\_t is a
genuine intrinsic E1 potential. Constants can be subtracted to normalize
it; the objective is unchanged because nu is a probability measure and e
translates by the same constant. By maximality and q\_t\textless=u+t f,

\[0\geq e(q_t)-e(u)-\int(q_t-u)\,d\nu
\geq e(q_t)-e(u)-t\int f\,d\nu.\]

For t\textgreater0 divide by t; for t\textless0 reverse the inequality.
The two outer derivatives at this actual unbounded E1 base were
reconstructed in134 and equal integral f dmu\_u. Taking the two limits
yields equality with integral f dnu for every original continuous
f.~Radon uniqueness proves

\[M_\theta(u)=V\nu\]

on the original entire X\^{}an. This is derived from the preconstructed
raw metric operator; no inverse-lift operator is defined. No PL-only or
quotient-carrier equation substitutes for it.

\subsection{5. Separation and uniqueness
everywhere}\label{separation-and-uniqueness-everywhere}

In particular every delta\_v at an original divisorial v has some
intrinsic solution rho\_v with mu\_rho\_v=delta\_v. This conclusion only
needed Section2's linear growth and the genuine
compactness/usc/derivative argument, not an ample-plus-divisor
decomposition. The solution rho\_v need not be bounded or PL.

Let u,u' be normalized intrinsic potentials with probability-normalized
numerical I\_theta(u,u')=0. The transferred four-point Holder estimate
of134, whose bounded theta proof is the epsilon=0 limit of the complete
BJ Lemma7.30 and whose intrinsic extension passes actual mixed weighted
floors, gives

\[\left|\int(u-u')\,(d\mu_{\rho_v}-d\mu_0)\right|
\leq C_n I_\theta(u,u')^{\alpha_n} I_\theta(\rho_v,0)^{1/2}
\max\{J(u),J(u'),J(\rho_v),J(0)\}^{1/2-\alpha_n},\]

with alpha\_n=2\^{}\{-n\}; harmless zero cases are interpreted by the
original finite estimate, or by adding one to the maximum. All factors
except the zero power of I are finite. The left side is
\textbar u(v)-u'(v)\textbar{} because mu\_0=delta\_v\_triv and both
generic values are zero. Thus u and u' agree at EVERY divisorial point.
Embed both in one actual ample Omega and use that theory's all-point
divisorial determination. They agree on the whole original X\^{}an,
including poles and nondivisorial points.

Conversely equal normalized functions have zero I by its actual weighted
formula. If two intrinsic potentials have the same raw metric measure,
that formula directly gives I=0; normalize and apply the preceding
argument. Therefore the original intrinsic nef operator is injective
modulo constants whenever V\textgreater0, provided133--135 have been
reconstructed as stated. Every atomic solution in Sections3--4 is unique
modulo constants. This also supplies uniqueness for ANY other target for
which an intrinsic solution is separately shown to exist.

\subsection{6. Exact boundary and primary
applications}\label{exact-boundary-and-primary-applications}

This proves existence and uniqueness for every finite divisorial atomic
probability target, and all-point separation of intrinsic I for every
positive-mass nef actual L. General finite-energy targets still need
attainment: their target integral has not yet been proved weakly
continuous on energy sublevels. A finite-coordinate argument cannot be
silently applied to such a target.

The whole actual-ample G0 theorem is the sole equation grant and
supplies only the retained ample embedding, determination and
linear-growth bounds. The additional primary statement/proof units used
here are BJ v4 Proposition1.28 and Lemma4.46 with the complete
coordinates listed in Section2, and Lemma7.30 with its original AMP
domain retained before polynomial transfer. The derivative and
compactness constructions are exactly the unreviewed proposals133--135,
already frozen for independent124; this new consequence was not fed to
that active review. Original P16.5 remains open and no new PASS is
asserted.

\section{Appendix E. Intrinsic target-integral continuity and attained
nef-and-big
equations}\label{appendix-e.-intrinsic-target-integral-continuity-and-attained-nef-and-big-equations}

The following complete mathematical derivation is retained as supporting
proof, with its earlier scope statements. It is not a historical-verdict
grant. Equation numbers in this appendix are local to it. The main
candidate and the fresh review criteria govern the new combined scope.

\section{A proposed full intrinsic variational theorem for nef and big
bundles}\label{a-proposed-full-intrinsic-variational-theorem-for-nef-and-big-bundles}

This is an UNREVIEWED own reconstruction. Its theorem concerns the
entire intrinsic E1 space of an actual nef and big rational bundle, with
canonical zero reference, positive raw top mass, and the original
analytic carrier. It still does not solve original P16.5 for general
nonnef big references or nonbig nonzero numerical classes.
Proposals132--135 and the corrected polar proposal136-repair2 must be
independently reconstructed; they are not equation grants. Neither
frozen native123 nor native124 is changed. The exact actual-ample G0
theorem is the only equation grant.

\subsection{1. Domain and the proposed
conclusion}\label{domain-and-the-proposed-conclusion}

Let k be algebraically closed, trivially valued, characteristic zero and
of arbitrary cardinality. Let X be smooth integral projective of
dimension n\textgreater=1, and L an actual nef and numerically big
rational line bundle with V=L\^{}n\textgreater0. Put theta=c1(L). Retain
the entire genuine numerical cone PSH(theta), its canonical zero metric
and the entire intrinsic floor-defined E1(theta). Construct the original
raw operator before targets, as in133--134. Write e=E\_theta/V,
mu\_u=M\_theta(u)/V, and C(u)=-e(u) for normalized u(v\_triv)=0.

The proposed theorem is: every intrinsic finite-dual-energy Radon
probability target nu admits a unique normalized genuine u in E1(theta)
with M\_theta(u)=V nu. Both original continuous-direction outer
derivatives, arbitrary decreasing mixed/weighted nets for a fixed theta,
full floor mass and all-point uniqueness are retained. The dual space is
the actual intrinsic one: D\_theta(nu)=sup(e(u)-integral u
dnu)\textless infinity, equivalently all intrinsic E1 tests are
integrable by135. No restriction to an ambient unbounded pairing space
is made.

The additional original gates beyond132--135 are treated below: PL
genuine approximation without semiampleness, uniform PL-direction outer
estimates, weak approximation of every target by actual metric measures,
target-integral continuity, and attainment. Atomic attainment and
all-point separation have an independent short route in137.

\subsection{2. An actual ample-plus-divisor
decomposition}\label{an-actual-ample-plus-divisor-decomposition}

Choose an actual ample H. Numerical bigness means theta is in the
interior of the closed cone generated by effective divisors. For
sufficiently small rational epsilon\textgreater0, theta-epsilon c1(H) is
pseudoeffective. Approximate that class closely by a finite effective
rational combination D, so that theta-c1(D) remains ample; rational
approximation and openness of the ample cone suffice. Smoothness makes
the divisors Cartier. The actual rational bundle A=L-O(D) is numerically
ample and hence actual ample. Thus L=A+D with A ample and D effective
rational, on the same X and chosen before the target. This is the
complete elementary numerical reconstruction already included in the
reviewed scoped119 paper Section3.1; it does not assume section-ring
finite generation or a contraction.

Let sigma=log\textbar\textbar s\_D\textbar\textbar{} for the canonical
trivial divisor metric, Z=\textbar D\textbar\^{}an. Multiplication of
actual sections of A by the canonical section of D gives the generically
finite actual L-FS function sigma. Numerical genuine membership follows
from the primary class of generically finite FS functions and its
section approximation, as in119 Section3.2. In particular sigma is
theta-psh and normalized at the generic trivial valuation. If D=0,
sigma=0 and Z is empty.

The corrected136 polar construction gives an intrinsic E1 function
chi(sigma) with exactly its poles. Therefore every intrinsic
finite-dual-energy nu, and every intrinsic metric output, gives Z zero
mass. The argument requires no claim that sigma itself is E1(theta).

\subsection{3. Bounded PL approximation in the genuine nef
cone}\label{bounded-pl-approximation-in-the-genuine-nef-cone}

Fix a normalized bounded theta-psh u with -k\textless=u\textless=0, k a
positive integer. We construct bounded CONTINUOUS PL theta-psh functions
psi\_i which approach u in numerical I\_theta and energy, and whose
integrals against every fixed Radon nu with nu(Z)=0 approach integral u
dnu. The target need not be selected when constructing this family.

For every rational epsilon\textgreater0, u is genuinely (theta+epsilon
c1(A))-psh. The granted fixed-ample approximation theorem gives finite
continuous FS majorants. Their finitely many real weights can be rounded
upward to rational weights with arbitrarily small uniform error;
adjoining the positive rational shifts below therefore gives rational FS
majorants as well. As in the explicit nef intersection proof121, add
positive rational shifts, order by epsilon decreasing, shift decreasing
and function decreasing, and use compactness of X\^{}an to compare
finite continuous old majorants. This yields a decreasing net g\_i of
finite continuous FS functions of actual L+epsilon\_i A, with epsilon\_i
decreasing to0, g\_i decreasing to u, and g\_i(v\_triv)=c\_i decreasing
to0. Take a tail c\_i\textless=1. Clamp at -k if needed, which remains
finite FS in the actual ample class. Thus
-k\textless=g\_i\textless=c\_i\textless=1.

Normalize and clamp once more:

\[\phi_i=\max(g_i-c_i,-k),\qquad -k\leq\phi_i\leq0.\]

These are rational finite continuous FS functions of L+epsilon\_i A.
They need not be a decreasing net, but
g\_i-c\_i\textless=phi\_i\textless=g\_i, so their uniform differences
from g\_i are at most c\_i. In every common ample class containing these
classes, energy monotonicity between these two shifts proves energy
convergence, while their divisorial values approach those of u. The
granted ample strong-topology criterion therefore gives strong
convergence there. This application is only to bounded inputs in that
actual ample domain.

Set a\_i=epsilon\_i/(1+epsilon\_i) and form the ACTUAL generically
finite L-FS function

\[\eta_i=\frac{\phi_i+\epsilon_i\sigma}{1+\epsilon_i}.\]

The identity of actual bundles is

\[(L+\epsilon_i A)+\epsilon_i D=(1+\epsilon_i)L.\]

Clearing rational powers and multiplying canonical divisor sections
therefore constructs eta\_i by actual L sections, with its poles on Z.
Define

\[\psi_i=\max(\eta_i,-k).\tag{1}\]

The constant -k is genuinely theta-psh because L is nef. The genuine
maximum rule makes psi\_i theta-psh. It is a finite CONTINUOUS PL
function: after clearing powers, eta\_i is the maximum of finitely many
rational weighted logarithms of coherent section ideals; adding the
constant branch -k gives the flag-ideal PL function. We do not claim
psi\_i is a finite fixed-L FS function, which would incorrectly imply
semiampleness. Its actual model metric and theta-psh membership are the
relevant facts. It is normalized, bounded between -k and0, and hence
belongs to the published ambient pairing domain on its restriction to
X\^{}lin.

\subsubsection{3.1. Energy and numerical I
convergence}\label{energy-and-numerical-i-convergence}

Put r\_i=max(a\_i sigma,-k). This function is genuinely a\_i theta-psh,
hence lies in the common actual ample class Omega=2theta+c1(A). As
epsilon decreases, r\_i increases and has zero limit at every divisorial
point, where sigma is finite. It is bounded, so the complete original
AMP increasing-net pairing theorem BJ7.38 gives E\_Omega(r\_i)
increasing to0. Its normalized generic values are0; the actual ample
strong-topology theorem gives r\_i strongly tending to0 in E1(Omega).

On a tail epsilon\_i\textless1, all phi\_i+r\_i also lie in Omega: the
difference between Omega and (1+a\_i)theta+epsilon\_i c1(A) is the ample
class (1-a\_i)theta+(1-epsilon\_i)c1(A). From
-k\textless=phi\_i\textless=0 and sigma\textless=0 we obtain the
pointwise squeeze

\[\phi_i+r_i\leq\psi_i\leq\phi_i+\epsilon_i k.\tag{2}\]

The elementary lower inequality uses max(x+y,-k)\textgreater=x+max(y,-k)
for -k\textless=x\textless=0, and eta\_i\textgreater=phi\_i+a\_i sigma.
Both outer functions are genuine Omega-psh. The energy of phi\_i tends
to that of u in Omega. Addition is continuous in the granted ample
strong topology, by the complete mixed Holder/addition proof
BJ7.34/7.36, applied in 2Omega if needed. Hence phi\_i+r\_i tends
strongly to u in 2Omega. Its Omega energy also tends to E\_Omega(u):
expand (Omega,phi\_i+r\_i)=(2Omega,phi\_i+r\_i)-(Omega,0), and apply the
published mixed pairing continuity to the verified ample factors. The
upper bound in (2) has the same Omega energy limit. Monotonicity then
gives E\_Omega(psi\_i) tending to E\_Omega(u). On every divisorial
point, psi\_i tends to u because sigma is finite there. Thus psi\_i
tends strongly to u in the one actual ample E1(Omega).

For theta energy and theta metric pairings, use the ACTUAL metric
difference

\[(\theta,v)=(\Omega,v)-(\Omega-\theta,0),\]

where Omega-theta=theta+c1(A) is ample. Every varying v here is bounded,
genuinely Omega-psh, and strongly convergent in that ample cone. Finite
multilinear expansion and the published mixed Holder theorem give

\[E_\theta(\psi_i)\to E_\theta(u),\qquad
\int u\,(d\mu_{\psi_i}-d\mu_u)\to0,\]

and similarly with psi\_i in the tested slot. Thus J\_u(psi\_i) tends
to0 and I\_theta(psi\_i,u) tends to0 by the bounded numerical
equivalence transferred in134. The auxiliary restriction to X\^{}lin
serves only to justify those published pairings; all functions and
measures remain on the whole X\^{}an.

\subsubsection{3.2. Target integrals and arbitrary-cardinality
nets}\label{target-integrals-and-arbitrary-cardinality-nets}

The decreasing CONTINUOUS bounded net g\_i has usc endpoint u, so
integral g\_i dnu decreases to integral u dnu for every Radon nu. This
specific assertion follows from compact Radon regularity: any bounded
usc function can be approximated from above in integral by continuous
majorants; directed compact covers by the open comparisons
g\_i\textless h+delta give a common later index below h+delta. It is not
monotone convergence for arbitrary Borel nets. Uniform differences c\_i
tending to0 give the same convergence for phi\_i.

From (2),

\[|\psi_i-\phi_i|\leq-r_i+\epsilon_i k.\]

The functions r\_i depend only on scalar epsilon\_i, are bounded, and
increase to0 outside Z, with value -k on Z. Since nu(Z)=0, their
integrals tend to0. This scalar convergence can be checked on
epsilon=1/j and then squeezed for the full scalar net; it uses ordinary
countable bounded convergence, not a general Borel-net assertion. Hence
integral psi\_i dnu tends to integral u dnu.

\subsubsection{3.3. Entire intrinsic
potentials}\label{entire-intrinsic-potentials}

For normalized u in E1(theta), its literal floors u\_k decrease to u,
have energies approaching E\_theta(u), and I\_theta(u\_k,u) tending to0
by134. All finite-dual-energy nu integrate u, so the COUNTABLE floors
also converge in nu integral. Approximate each fixed floor by the net in
Sections3.1--3.2. For each prescribed positive energy, I and integral
tolerance choose k and then a sufficiently late bounded PL approximant.
This supplies genuine bounded PL approximants to u with bounded energy
costs, numerical I convergence and nu-integral convergence. A single
countable sequence for a fixed target can be chosen from these
tolerances; this does not claim pointwise sequential approximation of
all psh functions on a nonmetrizable space. Consequently D\_theta(nu)
admits a maximizing net entirely in PL(X) intersect PSH(theta).

\subsection{4. A uniform original PL-direction
estimate}\label{a-uniform-original-pl-direction-estimate}

Here is the key bounded estimate, proved before its intrinsic extension.
For bounded genuine theta-psh q,u and any bounded genuine theta-psh v,
integration by parts and telescoping yield

\[\left|\int v\,(d\mu_q-d\mu_u)\right|\leq2n\|q-u\|_\infty.\tag{3}\]

Indeed each telescoping integral is the integral of q-u against the
difference of two positive mixed theta-products, one with theta\_v and
one with theta\_0 in the remaining slot. Each has probability mass1, so
its absolute value is at most twice the sup norm; there are n slots.
This is an actual bounded metric identity, also obtainable as the
epsilon=0 polynomial limit of the original AMP Lemma7.32. Constants in v
cancel. For intrinsic q,u at finite sup distance, pass their literal
floors in134's weighted convergence with fixed bounded test v. The same
estimate holds, with no dependence on their energies or v's energy.

For arbitrary original PL f, Corollary2.26 gives a rational
lambda\textgreater0 and bounded continuous A-psh PL functions b,b' with
b-b'=lambda f.~Split a rational FS linear combination into positive and
negative sums, divide by a rational total-weight bound, and add the zero
A metric for unused weight. This is a difference representation, not a
false claim that both signs of a small PL function are A-psh.

Let v\_k=max(sigma+b,-k), v'\_k=max(sigma+b',-k). Both are bounded
genuine theta-psh by actual section multiplication and the nef maximum
rule. Their difference d\_k is bounded in absolute value by
lambda\textbar\textbar f\textbar\textbar, and equals lambda f off a deep
sigma collar. Every intrinsic output gives Z zero mass by136-repair2.
For each fixed pair q,u, ordinary COUNTABLE bounded convergence as k
increases gives

\[\int d_k\,(d\mu_q-d\mu_u)\to
\lambda\int f\,(d\mu_q-d\mu_u).\]

Apply (3) separately to v\_k,v'\_k, whose bound is independent of k. We
obtain

\[\left|\int f\,(d\mu_q-d\mu_u)\right|
\leq\frac{4n}{\lambda}\|q-u\|_\infty.\tag{4}\]

This is uniform over the ENTIRE intrinsic E1 cone. It does not require
general unbounded sigma integrability or ambient E\textasciitilde1
membership.

For u in E1(theta), q\_s=P\_theta(u+s f) satisfies
\textbar\textbar q\_s-u\textbar\textbar\textless=\textbar s\textbar\textbar\textbar f\textbar\textbar.
The supporting inequalities proved in132 for continuous obstacles pass
to continuous majorants h\_i decreasing to u, with q\_i=P\_theta h\_i
decreasing to u and P\_theta(h\_i+s f) decreasing to q\_s. Entire
weighted weak convergence and energy convergence are134. At fixed s this
yields the two quotient bounds by integral f dmu\_u and integral f
dmu\_q\_s, with their order reversed for negative s. Combining with (4)
gives

\[\left|e(q_s)-e(u)-s\int f\,d\mu_u\right|
\leq K_f s^2,\qquad K_f=4n\|f\|_\infty/\lambda.\tag{5}\]

This is the uniform error needed for maximizing nets. We do not infer
uniform differentiability merely from pointwise existence of both
derivatives.

\subsection{5. Maximizing measures approach the actual
target}\label{maximizing-measures-approach-the-actual-target}

Let u\_i be any intrinsic maximizing net for a finite-dual-energy nu,
normalized at the generic point, and let
delta\_i=D\_theta(nu)-e(u\_i)+integral u\_i dnu tend to0. Coercivity135
bounds C(u\_i) eventually. The pointwise envelope q=P\_theta(u\_i+s f)
is below u\_i+s f and has finite energy. Its objective is at most
D\_theta(nu), so

\[e(q)-e(u_i)-s\int f\,d\nu\leq\delta_i.\]

By (5), for s\textgreater0,

\[\int f\,(d\mu_{u_i}-d\nu)\leq\delta_i/s+K_f s.\]

Replace f by -f to obtain the other bound. First let i tend to infinity
at fixed s, then s decrease to0. Thus mu\_u\_i tends weakly to nu
against EVERY original PL f.~Uniform PL density and the common unit
masses give weak convergence against the full C0(X\^{}an). This
conclusion is on the original carrier, not a restricted observable
family or an inverse definition.

\subsection{6. Uniform comparison with a target
integral}\label{uniform-comparison-with-a-target-integral}

Define J\_nu(psi)=D\_theta(nu)-e(psi)+integral psi dnu. Strict numerical
concavity of e, transferred in134, gives for intrinsic phi,psi

\[I_\theta(\phi,\psi)\leq C_n\bigl(J_\nu(\phi)+J_\nu(\psi)\bigr).\tag{6}\]

To see this directly, test the dual supremum at (phi+psi)/2. Its target
integral is the average of the two integrals, while the strict-concavity
energy excess bounds I below. All quantities are finite by the actual
intrinsic target characterization. This is the original Lemma9.15
argument with only transferred numerical identities.

Choose, using Section3.3, a maximizing net rho\_j in bounded genuine
PL(theta). Coercivity bounds C(rho\_j) eventually. For an intrinsic psi
with C(psi)\textless=B, and a bounded genuine PL(theta) test phi with
C(phi)\textless=B, the four-point Holder estimate of134 gives

\[\left|\int\phi\,(d\mu_{\rho_j}-d\mu_\psi)\right|
\leq K_B I_\theta(\rho_j,\psi)^{1/2}
\leq K'_B\bigl(J_\nu(\rho_j)+J_\nu(\psi)\bigr)^{1/2},\]

where B can be enlarged to include D\_theta(nu) and all eventual costs.
The fixed phi is CONTINUOUS, so Section5's original weak convergence
passes its integral to nu. We obtain

\[\left|\int\phi\,(d\nu-d\mu_\psi)\right|
\leq K'_B J_\nu(\psi)^{1/2}.\tag{7}\]

For arbitrary intrinsic phi of cost at most B, use the approximants in
Section3.3 for this fixed target nu. Their costs stay bounded, their nu
integrals approach that of phi, and their numerical I distance to phi
tends to0. The transferred Holder estimate in the TEST variable gives
convergence of their integrals against the fixed mu\_psi. Passing this
approximation in (7) gives the same estimate for every intrinsic phi,
uniformly on the indicated cost sublevel. No weak convergence theorem
with a discontinuous test was assumed in this step.

\subsection{7. Target integration is weakly continuous on cost
sublevels}\label{target-integration-is-weakly-continuous-on-cost-sublevels}

Because each rho\_j is an actual bounded PL model metric, the
independently constructed operator of132/primary Proposition7.19(ii) is
a FINITE sum of original divisorial component Dirac measures, including
the raw central-fiber multiplicities. Its positivity and raw mass are
already verified. Thus the map

\[\phi\longmapsto\int\phi\,d\mu_{\rho_j}\]

is a finite sum of the coordinates defining the weak topology of the
entire normalized genuine theta cone. In particular it is continuous
there. Equation(7), with psi=rho\_j and J\_nu(rho\_j) tending to0,
proves uniform convergence of these continuous functionals to phi
mapping to integral phi dnu on every cost sublevel C(phi)\textless=B.
Therefore the latter target integral is continuous in that relative weak
topology. This replaces the invalid shortcut of applying the ample
Proposition9.19 directly to the larger unbounded intrinsic nef domain.

The proof needs actual PL theta metrics; their finite support was proved
via the model formula, not inferred from a formal probability operator.
The numerical reference and component coefficients were retained through
the original construction.

\subsection{8. Attainment, equation and
uniqueness}\label{attainment-equation-and-uniqueness}

Choose a normalized maximizing sequence for nu. Coercivity gives an
eventual cost bound B, and the genuine nef compactness and intrinsic
weak energy usc in135 give a weakly convergent subnet with genuine
endpoint u in the same compact intrinsic cost sublevel. Section7 makes
its nu integral converge. Energy usc therefore implies

\[e(u)-\int u\,d\nu\geq\limsup_i\left(e(u_i)-\int u_i\,d\nu\right)=D_\theta(\nu).\]

The endpoint attains the supremum. Apply the two original outer
derivative signs of134 to q\_t=P\_theta(u+t f) for EVERY original f in
C0(X\^{}an), exactly as in137 Section4. Positivity and the obstacle
inequality give both inequalities for the target integral; the
derivatives imply equality with integral f dmu\_u. Hence M\_theta(u)=V
nu as original Radon measures.

All-point uniqueness follows from137: every divisorial Dirac target has
a solution by the direct finite-coordinate compactness argument,
independent of the general Section7. Four-point Holder against that
solution and the canonical zero measure proves that I=0 separates all
normalized genuine potentials at every divisorial point. The accepted
ample all-point determination then separates them on the entire analytic
space. Equal metric measures imply I=0 by its actual weighted formula.
Thus all solutions are unique modulo constants, not only on a generic
complement, a skeleton or an observable quotient.

\subsection{9. Dependencies, independent reconstruction and the original
boundary}\label{dependencies-independent-reconstruction-and-the-original-boundary}

The auxiliary A and D supply approximants and estimates only. The raw
operator, intrinsic energy, genuine cone and target dual were
constructed beforehand and contain no such auxiliary choice. The
resulting equation and uniqueness therefore use the given L, canonical
reference and raw top mass V. The numerical-big decomposition in
Section2 is the elementary actual bundle argument already reviewed
in119; section-volume equality with L\^{}n is not separately
reconstructed here. It must be checked wherever a larger original
formulation requires that explicit volume identification.

This proposed theorem uses the full unreviewed intrinsic
floor/mixed/weighted/derivative reconstruction133--135, the corrected
polar proof136-repair2, and the independent atomic/separation
deduction137. They must all be checked together. A proof of only the new
formal variational steps is insufficient. The general nonnef
minimal/model-reference compatibility and nonbig nonzero-carrier
equations remain unresolved; even an accepted nef-and-big theorem would
not close the entire original P16.5.

Additional complete primary statement/proof units actually read while
constructing this note: BJv4 Definition2.1/Theorem2.2/Lemma2.3
(LF943--1005); Definition2.18/Proposition2.19 (LF1182--1240);
Proposition2.25 and Corollary2.26 (LF1313--1337); complete
Holder/addition proofs7.34/7.36 (LF4340--4446, with Corollary7.35 and
Lemma7.37 complete proofs additionally read at LF4447--4491); complete
increasing pairing theorem7.38 (LF4493--4526); complete original
variational proofs9.11/9.15/9.17/9.19 (LF5136--5262); complete
linear-growth proofs1.28/4.46 (LF922--931 and2703--2712); and the
complete original ample strong-topology/attainment proofs12.4/12.8
within LF6027--6146. Their AMP domains are retained before the displayed
bounded/coefficient/floor deductions; none grants the proposed nonample
equation. Primary URL: https://arxiv.org/pdf/1801.08229v4 . No
independent PASS or original closure is asserted for this draft.

\end{document}
